\documentclass[UTF8, 12pt, fontset=fandol, oneside]{ctexart}
\usepackage{researchnotes}

\title{凸优化：等式约束优化}
\author{}
\date{}

\begin{document}

\maketitle

\section*{10\quad 等式约束优化}

\subsection*{10.1\quad 等式约束优化问题}

本章讨论下述等式约束凸优化问题的求解方法：
\begin{equation}
\begin{aligned}
\text{minimize}\quad & f(x)\\
\text{subject to}\quad & Ax=b,
\end{aligned}
\tag{10.1}
\end{equation}
其中 $f:\mathbb{R}^n\to\mathbb{R}$ 是二次连续可微凸函数，$A\in\mathbb{R}^{p\times n}$，$\operatorname{rank}A=p<n$。对 $A$ 的假设意味着等式约束数少于变量数，并且等式约束互相独立。我们假定存在一个最优解 $x^\star$，并用 $p^\star$ 表示其最优值，
\[
p^\star=\inf\{f(x)\mid Ax=b\}=f(x^\star).
\]
之前曾说明过（见 \S4.2.3 或 \S5.5.3），点 $x^\star\in\operatorname{dom} f$ 是优化问题 (10.1) 的最优解的充要条件是，存在 $\nu^\star\in\mathbb{R}^p$ 满足
\begin{equation}
Ax^\star=b,\qquad \nabla f(x^\star)+A^T\nu^\star=0.
\tag{10.2}
\end{equation}
因此，求解等式约束优化问题 (10.1) 等价于确定 KKT 方程 (10.2) 的解，这是含有 $n+p$ 个变量 $x^\star,\nu^\star$ 的 $n+p$ 个方程的求解问题。第一组方程，$Ax^\star=b$，称为原可行方程，是线性的。第二组方程，$\nabla f(x^\star)+A^T\nu^\star=0$，称为对偶可行方程，通常是非线性的。同无约束优化一样，很少情况下可以用解析方法求解这些最优性条件。最重要的特殊情况是 $f$ 为二次函数，我们将在 \S10.1.1 分析这种问题。

任何等式约束优化问题都可以通过消除等式约束转化为等价的无约束问题，此后就可以用第 9 章的方法求解。另一种处理方法是利用无约束优化方法求解对偶问题（假设对偶函数是二次可微的），然后从对偶解中复原等式约束问题 (10.1) 的解。在 \S10.1.2 和 \S10.1.3 中我们将分别对消除方法和对偶方法进行简要的讨论。

本章主要内容是对 Newton 方法进行扩展，使之能够直接处理等式约束。很多情况下这些方法比将等式约束问题转换为无约束问题的方法更好。一个原因是消除方法（或对偶方法）通常会破坏问题的结构，比如稀疏性；与此相反，直接处理等式约束的方法能够利用问题的结构。另一个原因是概念上的：直接处理等式约束的方法可以被视为直接求解最优性条件 (10.2) 的方法。

\subsection*{10.1.1\quad 等式约束凸二次规划}

考虑等式约束凸二次规划问题
\begin{equation}
\begin{aligned}
\text{minimize}\quad & f(x)=(1/2)x^TPx+q^Tx+r\\
\text{subject to}\quad & Ax=b,
\end{aligned}
\tag{10.3}
\end{equation}
其中 $P\in\mathbb{S}_+^n$，$A\in\mathbb{R}^{p\times n}$。该问题不仅本身具有重要性，同时也是扩展 Newton 方法处理等式约束问题的基础。此时最优性条件 (10.2) 成为
\[
Ax^\star=b,\qquad Px^\star+q+A^T\nu^\star=0,
\]
我们可以将其写为
\begin{equation}
\begin{bmatrix}
P & A^T\\
A & 0
\end{bmatrix}
\begin{bmatrix}
x^\star\\
\nu^\star
\end{bmatrix}
=
\begin{bmatrix}
-q\\
b
\end{bmatrix}.
\tag{10.4}
\end{equation}
这组 $n+p$ 个变量 $x^\star,\nu^\star$ 的 $n+p$ 个线性方程称为等式约束优化问题 (10.3) 的 KKT 系统。系数矩阵称为 KKT 矩阵。

如果 KKT 矩阵非奇异，存在唯一最优的原对偶对 $(x^\star,\nu^\star)$。如果 KKT 矩阵奇异，但 KKT 系统有解，任何解都构成最优对 $(x^\star,\nu^\star)$。如果 KKT 系统无解，二次优化问题或者无下界或者无解。确实，在这种情况下存在 $v\in\mathbb{R}^n$ 和 $w\in\mathbb{R}^p$ 满足
\[
Pv+A^Tw=0,\qquad Av=0,\qquad -q^Tv+b^Tw>0.
\]
令 $\hat{x}$ 为任意可行点。对于任何 $t$，点 $x=\hat{x}+tv$ 都是可行解，并且
\[
\begin{aligned}
f(\hat{x}+tv)
&=f(\hat{x})+t(v^TP\hat{x}+q^Tv)+(1/2)t^2v^TPv\\
&=f(\hat{x})+t(-\hat{x}^TA^Tw+q^Tv)-(1/2)t^2w^TAv\\
&=f(\hat{x})+t(-b^Tw+q^Tv),
\end{aligned}
\]
它可以随着 $t\to\infty$ 一直减小没有下界。

\noindent\textbf{KKT 矩阵的非奇异性}\quad
之前曾假定 $P\in\mathbb{S}_+^n$，$\operatorname{rank}A=p<n$。在该假定下对 KKT 矩阵的非奇异性有若干等价条件：
\[
\mathcal{N}(P)\cap\mathcal{N}(A)=\{0\},
\]
即 $P$ 和 $A$ 没有共同的非平凡零空间；
\[
Ax=0,\ x\neq 0\Longrightarrow x^TPx>0,
\]
即 $P$ 在 $A$ 的零空间是正定的；
\[
F^TPF\succ 0,
\]
其中 $F\in\mathbb{R}^{n\times(n-p)}$ 是满足 $\mathcal{R}(F)=\mathcal{N}(A)$ 的矩阵。（见习题 10.1。）作为一个重要的特殊情况，我们指出如果 $P\succ 0$，KKT 矩阵必然非奇异。

\subsection*{10.1.2\quad 消除等式约束}

求解等式约束问题 (10.1) 的一种一般性方法是先消除等式约束，像 \S4.2.4 描述的那样，然后用无约束优化方法求解相应的无约束问题。我们首先确定矩阵 $F\in\mathbb{R}^{n\times(n-p)}$ 和向量 $\hat{x}\in\mathbb{R}^n$，用以参数化（仿射）可行集
\[
\{x\mid Ax=b\}=\{Fz+\hat{x}\mid z\in\mathbb{R}^{n-p}\}.
\]
这里 $\hat{x}$ 可以选用 $Ax=b$ 的任何特殊解，$F\in\mathbb{R}^{n\times(n-p)}$ 是值域为 $A$ 的零空间的任何矩阵。然后形成简化或消除等式约束后的优化问题
\begin{equation}
\text{minimize}\quad \tilde{f}(z)=f(Fz+\hat{x}),
\tag{10.5}
\end{equation}
这是变量 $z\in\mathbb{R}^{n-p}$ 的无约束问题。利用它的解 $z^\star$ 可以确定等式约束问题的解
\[
x^\star=Fz^\star+\hat{x}.
\]
我们也可以为等式约束问题构造一个最优的对偶变量 $\nu^\star$，这就是
\[
\nu^\star=-(AA^T)^{-1}A\nabla f(x^\star).
\]
为了说明这个表达式的正确性，我们必须验证对偶可行性条件
\begin{equation}
\nabla f(x^\star)+A^T\left(-(AA^T)^{-1}A\nabla f(x^\star)\right)=0
\tag{10.6}
\end{equation}
成立。为说明这一点，我们指出
\[
\begin{bmatrix}
F^T\\
A
\end{bmatrix}
\left(\nabla f(x^\star)-A^T(AA^T)^{-1}A\nabla f(x^\star)\right)=0,
\]
其中上面的矩阵块利用了 $F^T\nabla f(x^\star)=\nabla\tilde{f}(z^\star)=0$ 和 $AF=0$。既然左边的矩阵是非奇异的，该式意味着式 (10.6)。

\noindent\textbf{例 10.1}\quad
受资源约束的最优配置。我们考虑问题
\[
\begin{aligned}
\text{minimize}\quad & \sum_{i=1}^n f_i(x_i)\\
\text{subject to}\quad & \sum_{i=1}^n x_i=b,
\end{aligned}
\]
其中 $f_i:\mathbb{R}\to\mathbb{R}$ 是二次可微凸函数，$b\in\mathbb{R}$ 是问题的参数。我们将其解释为一种总量为 $b$（预算值）的资源在 $n$ 个相互独立的活动中进行最优配置的问题。我们可以利用下式消除（比如）$x_n$，
\[
x_n=b-x_1-\cdots-x_{n-1},
\]
它对应于下述选择
\[
\hat{x}=be_n,\qquad
F=
\begin{bmatrix}
I\\
-\mathbf{1}^T
\end{bmatrix}
\in\mathbb{R}^{n\times(n-1)}.
\]
于是简化问题为
\[
\text{minimize}\quad f_n(b-x_1-\cdots-x_{n-1})+\sum_{i=1}^{n-1}f_i(x_i),
\]
变量为 $x_1,\cdots,x_{n-1}$。

\noindent\textbf{选择消除矩阵}\quad
显然，$\mathbb{R}^{n\times(n-p)}$ 中任何满足 $\mathcal{R}(F)=\mathcal{N}(A)$ 的矩阵都可以用作消除矩阵，因此对矩阵 $F$ 存在很多可能的选择。如果 $F$ 是这样一个矩阵，并且 $T\in\mathbb{R}^{(n-p)\times(n-p)}$ 是非奇异的，那么 $\tilde{F}=FT$ 也是一个合适的消除矩阵，因为
\[
\mathcal{R}(\tilde{F})=\mathcal{R}(F)=\mathcal{N}(A).
\]
反之，如果 $F$ 和 $\tilde{F}$ 是任意两个合适的消除矩阵，那么总存在某个非奇异矩阵 $T$ 使得 $\tilde{F}=FT$。

如果使用 $F$ 消除等式约束，我们将求解无约束问题
\[
\text{minimize}\quad f(Fz+\hat{x}).
\]
而如果使用 $\tilde{F}$，我们将求解无约束问题
\[
\text{minimize}\quad f(\tilde{F}\tilde{z}+\hat{x})=f(F(T\tilde{z})+\hat{x}),
\]
这个问题和上面的等价，就是简单地通过坐标变换 $z=T\tilde{z}$ 得到。换言之，改变消除矩阵可以视为对简化问题进行变量变换。

\subsection*{10.1.3\quad 用对偶方法求解等式约束问题}

求解优化问题 (10.1) 的另一个途径是先求解其对偶问题，然后复原最优的原变量 $x^\star$，在 \S5.5.5 中对此进行了描述。优化问题 (10.1) 的对偶函数是
\[
\begin{aligned}
g(\nu)
&=-b^T\nu+\inf_x(f(x)+\nu^TAx)\\
&=-b^T\nu-\sup_x\left(( -A^T\nu)^Tx-f(x)\right)\\
&=-b^T\nu-f^\ast(-A^T\nu),
\end{aligned}
\]
其中 $f^\ast$ 是 $f$ 的共轭，因此对偶问题是
\[
\text{maximize}\quad -b^T\nu-f^\ast(-A^T\nu).
\]
既然根据定义存在最优解，该问题是严格可行的，因此 Slater 条件成立。于是强对偶性成立，最优对偶目标可以达到，即存在 $\nu^\star$ 满足 $g(\nu^\star)=p^\star$。

如果对偶函数 $g$ 是二次可微的，那么在第 9 章描述的无约束优化方法可以用于极大化 $g$。（一般情况下，即使 $f$ 是二次可微的，对偶函数 $g$ 也不一定二次可微。）一旦找到最优的对偶变量 $\nu^\star$，我们就可以由它构造出原问题的最优解 $x^\star$。（并非总能简单完成；见 \S5.5.5。）

\noindent\textbf{例 10.2}\quad
等式约束的解析中心。我们考虑问题
\begin{equation}
\begin{aligned}
\text{minimize}\quad & f(x)=-\sum_{i=1}^n\log x_i\\
\text{subject to}\quad & Ax=b,
\end{aligned}
\tag{10.7}
\end{equation}
其中 $A\in\mathbb{R}^{p\times n}$，隐含约束 $x\succ 0$。利用
\[
f^\ast(y)=\sum_{i=1}^n(-1-\log(-y_i))=-n-\sum_{i=1}^n\log(-y_i),
\]
（$\operatorname{dom} f^\ast=-\mathbb{R}_{++}^n$），对偶问题为
\begin{equation}
\text{maximize}\quad g(\nu)=-b^T\nu+n+\sum_{i=1}^n\log(A^T\nu)_i,
\tag{10.8}
\end{equation}
隐含约束 $A^T\nu\succ 0$。这里我们可以很容易的求解对偶可行性方程，即确定极大化 $L(x,\nu)$ 的 $x$：
\[
\nabla f(x)+A^T\nu=-(1/x_1,\cdots,1/x_n)+A^T\nu=0,
\]
于是
\begin{equation}
x_i(\nu)=1/(A^T\nu)_i.
\tag{10.9}
\end{equation}
为求解等式约束的解析中心问题 (10.7)，我们先求解（无约束的）对偶问题 (10.8)，然后利用式 (10.9) 复原问题 (10.7) 的最优解。

\subsection*{10.2\quad 等式约束的 Newton 方法}

本节讨论能够处理等式约束的扩展 Newton 方法。除了以下两点不同，该方法和没有约束的 Newton 方法几乎一样：初始点必须可行（即满足 $x\in\operatorname{dom} f$，$Ax=b$），根据等式约束的需要对 Newton 方向的定义进行了适当的修改。具体地说，我们要保证 Newton 方向 $\Delta x_{\mathrm{nt}}$ 是可行方向，即 $A\Delta x_{\mathrm{nt}}=0$。

\subsection*{10.2.1\quad Newton 方向}

\noindent\textbf{基于二阶近似的定义}\quad
为导出等式约束问题
\[
\begin{aligned}
\text{minimize}\quad & f(x)\\
\text{subject to}\quad & Ax=b,
\end{aligned}
\]
在可行点 $x$ 处的 Newton 方向 $\Delta x_{\mathrm{nt}}$，我们将目标函数换成其在 $x$ 附近的二阶 Taylor 近似，形成下述问题
\begin{equation}
\begin{aligned}
\text{minimize}\quad & \hat{f}(x+v)=f(x)+\nabla f(x)^Tv+(1/2)v^T\nabla^2 f(x)v\\
\text{subject to}\quad & A(x+v)=b,
\end{aligned}
\tag{10.10}
\end{equation}
该问题的变量为 $v$。这是一个（凸的）带等式约束的二次极小问题，可以用解析方法求解。我们假定有关的 KKT 矩阵非奇异，在此基础上定义 $x$ 处的 Newton 方向 $\Delta x_{\mathrm{nt}}$ 为凸二次问题 (10.10) 的解。换言之，Newton 方向 $\Delta x_{\mathrm{nt}}$ 是一个向量，将其加到 $x$ 上就形成以 $f$ 的二阶近似为目标函数的对应问题的最优解。

根据 \S10.1.1 对等式约束二次问题的分析，Newton 方向 $\Delta x_{\mathrm{nt}}$ 由以下方程确定
\begin{equation}
\begin{bmatrix}
\nabla^2 f(x) & A^T\\
A & 0
\end{bmatrix}
\begin{bmatrix}
\Delta x_{\mathrm{nt}}\\
w
\end{bmatrix}
=
\begin{bmatrix}
-\nabla f(x)\\
0
\end{bmatrix},
\tag{10.11}
\end{equation}
其中 $w$ 是该二次问题的最优对偶变量。Newton 方向只在 KKT 矩阵非奇异的点有定义。

如同用 Newton 方法求解无约束问题一样，当目标函数 $f$ 是严格的二次函数时，Newton 修正向量 $x+\Delta x_{\mathrm{nt}}$ 是等式约束极小化问题的准确最优解，在这种情况下向量 $w$ 是原问题最优的对偶变量。由此想到，正如无约束情况一样，当 $f$ 接近二次时，$x+\Delta x_{\mathrm{nt}}$ 应该是最优解 $x^\star$ 的很好估计，而 $w$ 则应该是最优的对偶变量 $\nu^\star$ 的很好估计。

\noindent\textbf{线性化最优性条件的解}\quad
我们可以将 Newton 方向 $\Delta x_{\mathrm{nt}}$ 及其相关向量 $w$ 解释为最优性条件
\[
Ax^\star=b,\qquad \nabla f(x^\star)+A^T\nu^\star=0
\]
的线性近似方程组的解。我们用 $x+\Delta x_{\mathrm{nt}}$ 代替 $x^\star$，用 $w$ 代替 $\nu^\star$，并将第二个方程中的梯度项换成其在 $x$ 附近的线性近似，从而得到
\[
A(x+\Delta x_{\mathrm{nt}})=b,\qquad
\nabla f(x+\Delta x_{\mathrm{nt}})+A^Tw\approx
\nabla f(x)+\nabla^2 f(x)\Delta x_{\mathrm{nt}}+A^Tw=0.
\]
利用 $Ax=b$，以上方程变成
\[
A\Delta x_{\mathrm{nt}}=0,\qquad
\nabla^2 f(x)\Delta x_{\mathrm{nt}}+A^Tw=-\nabla f(x),
\]
这和定义 Newton 方向的方程 (10.11) 完全一样。

\noindent\textbf{Newton 减量}\quad
我们将等式约束问题的 Newton 减量定义为
\begin{equation}
\lambda(x)=\left(\Delta x_{\mathrm{nt}}^T\nabla^2 f(x)\Delta x_{\mathrm{nt}}\right)^{1/2}.
\tag{10.12}
\end{equation}
这和无约束情况的表示 (9.20) 完全一样，因此也可以进行同样的解释。例如，$\lambda(x)$ 是 Newton 方向的 Hessian 矩阵范数。

用
\[
\hat{f}(x+v)=f(x)+\nabla f(x)^Tv+(1/2)v^T\nabla^2 f(x)v
\]
表示 $f$ 在 $x$ 处的二阶 Taylor 近似。$f(x)$ 和二次模型之间的差值满足
\begin{equation}
f(x)-\inf\{\hat{f}(x+v)\mid A(x+v)=b\}=\lambda(x)^2/2,
\tag{10.13}
\end{equation}
同无约束情况完全一样（见习题 10.6）。这说明，如同无约束情况，$\lambda(x)^2/2$ 对 $x$ 处的 $f(x)-p^\star$ 给出了基于二次模型的一个估计，而 $\lambda(x)$（或者 $\lambda(x)^2$ 的倍数）也可以作为设计好的停止准则的基础。

Newton 减量也出现在直线搜索中，因为 $f$ 沿方向 $\Delta x_{\mathrm{nt}}$ 的方向导数是
\begin{equation}
\left.\frac{d}{dt}f(x+t\Delta x_{\mathrm{nt}})\right|_{t=0}
=\nabla f(x)^T\Delta x_{\mathrm{nt}}=-\lambda(x)^2,
\tag{10.14}
\end{equation}
和无约束情况一样。

\noindent\textbf{可行下降方法}\quad
假定 $Ax=b$。我们说 $v\in\mathbb{R}^n$ 是一个可行方向，如果 $Av=0$。在这种情况下，每个具有 $x+tv$ 形式的点也是可行解，即 $A(x+tv)=b$。我们说 $v$ 是 $f$ 在 $x$ 处的一个下降方向，如果对小的 $t>0$，有 $f(x+tv)<f(x)$。

Newton 方向总是可行下降方向（除非 $x$ 已经是最优解，此时 $\Delta x_{\mathrm{nt}}=0$）。确实，定义 $\Delta x_{\mathrm{nt}}$ 的第一组方程是 $A\Delta x_{\mathrm{nt}}=0$，说明它是可行方向；而式 (10.14) 则说明它也是下降方向。

\noindent\textbf{仿射不变性}\quad
同无约束优化问题一样，等式约束优化问题中的 Newton 方向和 Newton 减量也是仿射不变的。假定 $T\in\mathbb{R}^{n\times n}$ 是非奇异的，定义 $\bar{f}(y)=f(Ty)$。我们有
\[
\nabla \bar{f}(y)=T^T\nabla f(Ty),\qquad
\nabla^2\bar{f}(y)=T^T\nabla^2 f(Ty)T,
\]
而等式约束 $Ax=b$ 则成为 $ATy=b$。

现在考虑在 $ATy=b$ 的约束下极小化 $\bar{f}(y)$ 的问题。在 $y$ 处的 Newton 方向是 $\Delta y_{\mathrm{nt}}$，它是以下方程的解
\[
\begin{bmatrix}
T^T\nabla^2 f(Ty)T & T^TA^T\\
AT & 0
\end{bmatrix}
\begin{bmatrix}
\Delta y_{\mathrm{nt}}\\
\bar{w}
\end{bmatrix}
=
\begin{bmatrix}
-T^T\nabla f(Ty)\\
0
\end{bmatrix}.
\]
同式 (10.11) 给出的 $f$ 在 $x=Ty$ 处的 Newton 方向 $\Delta x_{\mathrm{nt}}$ 进行比较，可以看出
\[
T\Delta y_{\mathrm{nt}}=\Delta x_{\mathrm{nt}},
\]
（并且 $w=\bar{w}$，）即 $y$ 和 $x$ 处的 Newton 方向由和 $Ty=x$ 相同的坐标变换相联系。

\subsection*{10.2.2\quad 等式约束的 Newton 方法}

等式约束下 Newton 方法的框架和无约束情况完全一样。

\noindent\textbf{算法 10.1}\quad
等式约束优化问题的 Newton 方法。

给定满足 $Ax=b$ 的初始点 $x\in\operatorname{dom} f$，误差阈值 $\epsilon>0$。

\noindent\textbf{重复进行}

\noindent 1. 计算 Newton 方向和 Newton 减量 $\Delta x_{\mathrm{nt}},\lambda(x)$。

\noindent 2. 停止准则。如果 $\lambda^2/2\leq \epsilon$ 则退出。

\noindent 3. 直线搜索。通过回溯直线搜索确定步长 $t$。

\noindent 4. 修改。$x:=x+t\Delta x_{\mathrm{nt}}$。

这是一种可行下降方法，因为所有迭代点都是可行的，并且满足 $f(x^{(k+1)})<f(x^{(k)})$（除非 $x^{(k)}$ 已经是最优解）。Newton 方法需要每个 $x$ 处的 KKT 矩阵可逆；我们将在 \S10.2.4 对收敛性需要的假设进行更准确的描述。

\subsection*{10.2.3\quad Newton 方法和消除法}

我们现在说明，对等式约束问题 (10.1) 采用 Newton 方法的迭代过程，和对简化问题 (10.5) 采用 Newton 方法的迭代过程完全一致。假定 $F$ 满足 $\mathcal{R}(F)=\mathcal{N}(A)$ 和 $\operatorname{rank}F=n-p$，$\hat{x}$ 满足 $A\hat{x}=b$。简化目标函数 $\tilde{f}(z)=f(Fz+\hat{x})$ 的梯度和 Hessian 矩阵是
\[
\nabla\tilde{f}(z)=F^T\nabla f(Fz+\hat{x}),\qquad
\nabla^2\tilde{f}(z)=F^T\nabla^2 f(Fz+\hat{x})F.
\]
从 Hessian 矩阵可以看出，等式约束问题的 Newton 方向有定义，即 KKT 矩阵
\[
\begin{bmatrix}
\nabla^2 f(x) & A^T\\
A & 0
\end{bmatrix}
\]
可逆的充要条件是简化问题的 Newton 方向有定义，即 $\nabla^2\tilde{f}(z)$ 可逆。

简化问题的 Newton 方向是
\begin{equation}
\Delta z_{\mathrm{nt}}
=-\nabla^2\tilde{f}(z)^{-1}\nabla\tilde{f}(z)
=-(F^T\nabla^2 f(x)F)^{-1}F^T\nabla f(x),
\tag{10.15}
\end{equation}
其中 $x=Fz+\hat{x}$。简化问题的这个搜索方向对应于原始的等式约束问题的方向是
\[
F\Delta z_{\mathrm{nt}}=-F(F^T\nabla^2 f(x)F)^{-1}F^T\nabla f(x).
\]
我们将说明这个方向和式 (10.11) 定义的原始的等式约束问题的 Newton 方向 $\Delta x_{\mathrm{nt}}$ 完全一样。

为说明以上等价性，我们定义 $\Delta x_{\mathrm{nt}}=F\Delta z_{\mathrm{nt}}$，选择
\[
w=-(AA^T)^{-1}A\left(\nabla f(x)+\nabla^2 f(x)\Delta x_{\mathrm{nt}}\right),
\]
然后验证它们能满足定义 Newton 方向的方程
\begin{equation}
\nabla^2 f(x)\Delta x_{\mathrm{nt}}+A^Tw+\nabla f(x)=0,\qquad
A\Delta x_{\mathrm{nt}}=0.
\tag{10.16}
\end{equation}
第二个方程 $A\Delta x_{\mathrm{nt}}=0$ 成立是因为 $AF=0$。为了验证第一个方程，我们利用
\[
\begin{bmatrix}
F^T\\
A
\end{bmatrix}
\left(\nabla^2 f(x)\Delta x_{\mathrm{nt}}+A^Tw+\nabla f(x)\right)
=
\begin{bmatrix}
F^T\nabla^2 f(x)\Delta x_{\mathrm{nt}}+F^TA^Tw+F^T\nabla f(x)\\
A\nabla^2 f(x)\Delta x_{\mathrm{nt}}+AA^Tw+A\nabla f(x)
\end{bmatrix}
=0.
\]
因为第一行左边的矩阵是非奇异阵，我们可以断定式 (10.16) 成立。

类似地，下式说明 $\tilde{f}$ 在 $z$ 处的 Newton 减量 $\tilde{\lambda}(z)$ 和 $f$ 在 $x$ 处的 Newton 减量相等：
\[
\begin{aligned}
\tilde{\lambda}(z)^2
&=\Delta z_{\mathrm{nt}}^T\nabla^2\tilde{f}(z)\Delta z_{\mathrm{nt}}\\
&=\Delta z_{\mathrm{nt}}^TF^T\nabla^2 f(x)F\Delta z_{\mathrm{nt}}\\
&=\Delta x_{\mathrm{nt}}^T\nabla^2 f(x)\Delta x_{\mathrm{nt}}\\
&=\lambda(x)^2.
\end{aligned}
\]

\subsection*{10.2.4\quad 收敛性分析}

我们已经知道，用 Newton 方法求解等式约束问题和用 Newton 方法求解消除等式约束的简化问题完全等价。因此，我们关于无约束问题 Newton 方法收敛性的所有结果都可以推广于等式约束问题的 Newton 方法。特别是，用 Newton 方法求解等式约束问题的实际性能应该和用 Newton 方法求解无约束问题的性能完全一样。一旦 $x^{(k)}$ 接近 $x^\star$，收敛速度将非常快，经过几次迭代就可以获得非常高的精度。

\noindent\textbf{假设}\quad
我们给出以下假设。

\begin{itemize}
\item 下水平集
\[
S=\{x\mid x\in\operatorname{dom} f,\ f(x)\leq f(x^{(0)}),\ Ax=b\}
\]
是闭的，其中 $x^{(0)}\in\operatorname{dom} f$ 满足 $Ax^{(0)}=b$。该假设对闭的 $f$ 成立（见 \S A.3.3）。
\item 在集合 $S$ 上，我们有 $\nabla^2 f(x)\preceq MI$，
\begin{equation}
\left\|
\begin{bmatrix}
\nabla^2 f(x) & A^T\\
A & 0
\end{bmatrix}^{-1}
\right\|_2\leq K,
\tag{10.17}
\end{equation}
即 KKT 矩阵的逆矩阵在 $S$ 上有界。（当然，为了使 Newton 方向在 $S$ 的每个点有定义，该逆矩阵必须存在。）
\item 对于 $x,\tilde{x}\in S$，$\nabla^2 f$ 满足 Lipschitz 条件
\[
\|\nabla^2 f(x)-\nabla^2 f(\tilde{x})\|_2\leq L\|x-\tilde{x}\|_2.
\]
\end{itemize}

\noindent\textbf{KKT 矩阵逆有界假定}\quad
条件 (10.17) 的作用同分析标准 Newton 方法的强凸性假设类似（\S9.5.3，第 465 页）。如果没有等式约束，式 (10.17) 退化为在 $S$ 上 $\|\nabla^2 f(x)^{-1}\|_2\leq K$，因此我们可以当 $\nabla^2 f(x)\succeq mI$ 在 $S$ 上成立时取 $K=1/m$，其中 $m>0$。有等式约束时，这个条件不会像最小特征值存在正下界那样简单。因为 KKT 矩阵是对称的，条件 (10.17) 是指它的所有特征值，其中 $n$ 个为正，$p$ 个为负，和 $0$ 之间的距离都有下界。

\noindent\textbf{利用消除等式约束问题的分析}\quad
以上假设意味着消除等式约束后的目标函数 $\tilde{f}$，和相应的初始点 $z^{(0)}=\hat{x}+Fx^{(0)}$ 一起，满足 \S9.5.3 给出的无约束 Newton 方法收敛性分析的假定（其中参数 $\tilde{m}$，$\tilde{M}$ 和 $\tilde{L}$ 不同）。由此可知等式约束的 Newton 方法收敛于 $x^\star$（以及 $\nu^\star$）。

容易证明，上面的假设意味着消除等式约束后的问题满足无约束 Newton 方法的假定（见习题 10.4）。这里我们证明一个隐含的巧妙结果：KKT 矩阵逆有界的条件和上界假定 $\nabla^2 f(x)\preceq MI$ 一起，意味着 $\nabla^2\tilde{f}(z)\succeq mI$ 对某个正常数 $m$ 成立。更准确地说，我们要证明这个不等式对
\begin{equation}
m=\frac{\sigma_{\min}(F)^2}{K^2M}
\tag{10.18}
\end{equation}
成立，因为 $F$ 是满秩的，它确实是正的。

我们用反证法证明上述论断。假定 $F^THF\nsucceq mI$，其中 $H=\nabla^2 f(x)$。我们可以找到 $u$ 满足 $\|u\|_2=1$，$u^TF^THFu<m$，即 $\|H^{1/2}Fu\|_2<m^{1/2}$。利用 $AF=0$，我们有
\[
\begin{bmatrix}
H & A^T\\
A & 0
\end{bmatrix}
\begin{bmatrix}
Fu\\
0
\end{bmatrix}
=
\begin{bmatrix}
HFu\\
0
\end{bmatrix},
\]
因此
\[
\left\|
\begin{bmatrix}
H & A^T\\
A & 0
\end{bmatrix}^{-1}
\right\|_2
\geq
\frac{\left\|\begin{bmatrix}Fu\\0\end{bmatrix}\right\|_2}
{\left\|\begin{bmatrix}HFu\\0\end{bmatrix}\right\|_2}
=\frac{\|Fu\|_2}{\|HFu\|_2}.
\]
利用 $\|Fu\|_2\geq\sigma_{\min}(F)$ 和
\[
\|HFu\|_2\leq \|H^{1/2}\|_2\|H^{1/2}Fu\|_2<M^{1/2}m^{1/2},
\]
我们可断定
\[
\left\|
\begin{bmatrix}
H & A^T\\
A & 0
\end{bmatrix}^{-1}
\right\|_2
>\frac{\sigma_{\min}(F)}{M^{1/2}m^{1/2}}=K,
\]
其中用到我们在式 (10.18) 给出的 $m$ 的表示式。

\noindent\textbf{自和谐函数的收敛性分析}\quad
如果 $f$ 是自和谐的，那么 $\tilde{f}(z)=f(Fz+\hat{x})$ 同样也是。由此可知对于自和谐的 $f$，我们有和无约束问题完全一样的复杂性估计：产生满足精度 $\epsilon$ 的解所需要的迭代次数不会大于
\[
\frac{20-8\alpha}{\alpha\beta(1-2\alpha)^2}(f(x^{(0)})-p^\star)+\log_2\log_2(1/\epsilon),
\]
其中 $\alpha$ 和 $\beta$ 是回溯直线搜索的参数（见式 (9.56)）。

\section*{10.3\quad 不可行初始点的 Newton 方法}

如同 \S10.2 描述的那样，Newton 方法是一个可行下降方法。本节我们描述一种推广的 Newton 方法，它能够从不可行的初始点开始进行迭代。

\subsection*{10.3.1\quad 不可行点的 Newton 方向}

和 Newton 方法一样，我们从等式约束优化问题的最优性条件开始：
\[
Ax^\star=b,\qquad \nabla f(x^\star)+A^T\nu^\star=0.
\]
用 $x$ 表示当前点，我们不假定它是可行的，但假定它满足 $x\in\operatorname{dom} f$。我们的目的是找到一个方向 $\Delta x$ 使 $x+\Delta x$ 满足（至少近似满足）最优性条件，即 $x+\Delta x\approx x^\star$。为此我们在最优性条件中用 $x+\Delta x$ 代替 $x^\star$，用 $w$ 代替 $\nu^\star$，并利用梯度的一阶近似
\[
\nabla f(x+\Delta x)\approx \nabla f(x)+\nabla^2 f(x)\Delta x
\]
得到
\[
A(x+\Delta x)=b,\qquad
\nabla f(x)+\nabla^2 f(x)\Delta x+A^Tw=0.
\]
这是 $\Delta x$ 和 $w$ 的一组线性方程：
\begin{equation}
\begin{bmatrix}
\nabla^2 f(x) & A^T\\
A & 0
\end{bmatrix}
\begin{bmatrix}
\Delta x\\
w
\end{bmatrix}
=-
\begin{bmatrix}
\nabla f(x)\\
Ax-b
\end{bmatrix}.
\tag{10.19}
\end{equation}
这组方程和在可行点 $x$ 处定义 Newton 方向的式 (10.11) 相似，只有一点差别：右边第二块元素含有 $Ax-b$，它是线性等式约束的残差向量。如果 $x$ 是可行的，残差等于零，方程 (10.19) 退化为在可行点 $x$ 处定义 Newton 方向的方程 (10.11)。因此，如果 $x$ 是可行的，由式 (10.19) 定义的方向 $\Delta x$ 和前面描述的 Newton 方向（只对可行的 $x$ 有定义）一致。由于这个原因，我们用符号 $\Delta x_{\mathrm{nt}}$ 表示式 (10.19) 定义的方向 $\Delta x$，并在不会混淆时将它称为 $x$ 处的 Newton 方向。

\noindent\textbf{作为原对偶 Newton 方向的解释}\quad
我们可以基于等式约束问题的原对偶方法对方程 (10.19) 给出一种解释。所谓原对偶方法是指同时修改原变量 $x$ 和对偶变量 $\nu$ 使最优性条件（近似）满足的方法。

我们将最优性条件表示成 $r(x^\star,\nu^\star)=0$，其中 $r:\mathbb{R}^n\times\mathbb{R}^p\to\mathbb{R}^n\times\mathbb{R}^p$ 由下式定义
\[
r(x,\nu)=(r_{\mathrm{dual}}(x,\nu),r_{\mathrm{pri}}(x,\nu)).
\]
这里
\[
r_{\mathrm{dual}}(x,\nu)=\nabla f(x)+A^T\nu,\qquad
r_{\mathrm{pri}}(x,\nu)=Ax-b
\]
分别是对偶残差和原残差。在当前估计 $y$ 附近 $r$ 的一阶 Taylor 近似是
\[
r(y+z)\approx \hat{r}(y+z)=r(y)+Dr(y)z,
\]
其中 $Dr(y)\in\mathbb{R}^{(n+p)\times(n+p)}$ 是 $r$ 在 $y$ 处的导数（见 \S A.4.1）。我们将原对偶 Newton 方向 $\Delta y_{\mathrm{pd}}$ 定义为使 Taylor 近似 $\hat{r}(y+z)$ 等于 $0$ 的步径 $z$，即
\begin{equation}
Dr(y)\Delta y_{\mathrm{pd}}=-r(y).
\tag{10.20}
\end{equation}
注意此处我们将 $x$ 和 $\nu$ 都视为变量；$\Delta y_{\mathrm{pd}}=(\Delta x_{\mathrm{pd}},\Delta\nu_{\mathrm{pd}})$ 同时给出了原对偶问题的步径。

对 $r$ 求导数，可以将式 (10.20) 表示为
\begin{equation}
\begin{bmatrix}
\nabla^2 f(x) & A^T\\
A & 0
\end{bmatrix}
\begin{bmatrix}
\Delta x_{\mathrm{pd}}\\
\Delta\nu_{\mathrm{pd}}
\end{bmatrix}
=-
\begin{bmatrix}
r_{\mathrm{dual}}\\
r_{\mathrm{pri}}
\end{bmatrix}
=-
\begin{bmatrix}
\nabla f(x)+A^T\nu\\
Ax-b
\end{bmatrix}.
\tag{10.21}
\end{equation}
将 $\nu+\Delta\nu_{\mathrm{pd}}$ 写成 $\nu^+$，上式又可表示成
\begin{equation}
\begin{bmatrix}
\nabla^2 f(x) & A^T\\
A & 0
\end{bmatrix}
\begin{bmatrix}
\Delta x_{\mathrm{pd}}\\
\nu^+
\end{bmatrix}
=-
\begin{bmatrix}
\nabla f(x)\\
Ax-b
\end{bmatrix},
\tag{10.22}
\end{equation}
这和方程组 (10.19) 完全一样。因此，式 (10.19)、式 (10.21) 和式 (10.22) 的解之间存在以下联系
\[
\Delta x_{\mathrm{nt}}=\Delta x_{\mathrm{pd}},\qquad
w=\nu^+=\nu+\Delta\nu_{\mathrm{pd}}.
\]
这表明（不可行）Newton 方向和原对偶方向中的原问题对应向量相同，而相应的对偶向量 $w$ 则是修正的原对偶变量 $\nu^+=\nu+\Delta\nu_{\mathrm{pd}}$。

由式 (10.21) 和式 (10.22) 给出的 Newton 方向和对偶变量（或对偶方向）的两种表达式彼此等价，但每种表达式揭示了 Newton 方向一种不同的特点。方程 (10.21) 表明 Newton 方向和相应的对偶方向是以原对偶残差为右边项的方程组的解。而我们最初定义 Newton 方向的式 (10.22) 则给出了 Newton 方向和修正的对偶变量，从中可以看出对偶变量的当前值对计算对偶方向或者其修正值都不起作用。

\noindent\textbf{残差范数的缩减性质}\quad
在不可行点处的 Newton 方向不一定是 $f$ 的下降方向。从式 (10.19) 可以看出，
\[
\begin{aligned}
\left.\frac{d}{dt}f(x+t\Delta x)\right|_{t=0}
&=\nabla f(x)^T\Delta x\\
&=-\Delta x^T(\nabla^2 f(x)\Delta x+A^Tw)\\
&=-\Delta x^T\nabla^2 f(x)\Delta x+(Ax-b)^Tw,
\end{aligned}
\]
它并不一定是负的（当然，除非 $x$ 是可行的，即 $Ax=b$）。然而，原对偶解释表明残差范数沿 Newton 方向下降，即
\[
\left.\frac{d}{dt}\|r(y+t\Delta y_{\mathrm{pd}})\|_2^2\right|_{t=0}
=2r(y)^TDr(y)\Delta y_{\mathrm{pd}}
=-2r(y)^Tr(y).
\]
利用上面范数平方的导数，可以得到
\begin{equation}
\left.\frac{d}{dt}\|r(y+t\Delta y_{\mathrm{pd}})\|_2\right|_{t=0}
=-\|r(y)\|_2.
\tag{10.23}
\end{equation}
该式让我们能够利用 $\|r\|_2$ 检测不可行 Newton 方法的进展，例如，检测直线搜索的进展。（对于标准 Newton 方法，我们利用函数 $f$ 的取值检测算法的进展，至少在二次收敛阶段以前是这样。）

\noindent\textbf{完整步径的可行性质}\quad
由式 (10.19) 给出的 Newton 方向 $\Delta x_{\mathrm{nt}}$（根据定义）具有以下性质
\begin{equation}
A(x+\Delta x_{\mathrm{nt}})=b.
\tag{10.24}
\end{equation}
由此可知，如果沿 Newton 方向 $\Delta x_{\mathrm{nt}}$ 前进的步长等于 $1$，下一个迭代点将是可行的。一旦 $x$ 是可行点，Newton 方向就成为可行方向，不管以后的步长等于多少，所有的迭代点都将是可行点。

我们可以在更加一般的情况下分析阻尼步长对等式约束残差 $r_{\mathrm{pri}}$ 的影响。选定步长 $t\in[0,1]$，下一个迭代点是 $x^+=x+t\Delta x_{\mathrm{nt}}$，利用式 (10.24) 可求出其等式约束残差
\[
r_{\mathrm{pri}}^+=A(x+\Delta x_{\mathrm{nt}}t)-b=(1-t)(Ax-b)=(1-t)r_{\mathrm{pri}}.
\]
因此，阻尼步长 $t$ 能使残差缩减为原来的 $1-t$ 倍。现在假设对于 $i=0,\cdots,k-1$ 我们有
\[
x^{(i+1)}=x^{(i)}+t^{(i)}\Delta x_{\mathrm{nt}}^{(i)},
\]
其中 $\Delta x_{\mathrm{nt}}^{(i)}$ 是点 $x^{(i)}\in\operatorname{dom} f$ 处的 Newton 方向，$t^{(i)}\in[0,1]$。于是
\[
r^{(k)}=\left(\prod_{i=0}^{k-1}(1-t^{(i)})\right)r^{(0)},
\]
其中 $r^{(i)}=Ax^{(i)}-b$ 是 $x^{(i)}$ 的残差。该式表明迭代过程中原问题的残差向量始终和初始残差向量方向相同，但每次迭代后都按比例缩减。它也表明一旦某个步长等于 $1$，其后所有迭代点都是原问题的可行点。

\subsection*{10.3.2\quad 不可行初始点 Newton 方法}

利用式 (10.19) 定义的 Newton 方向 $\Delta x_{\mathrm{nt}}$，我们可以对 Newton 方法进行推广，使之能够处理 $x^{(0)}\in\operatorname{dom} f$，但不一定满足 $Ax^{(0)}=b$ 的情况。我们同时利用 Newton 方向的对偶部分：在式 (10.19) 中给出的 $\Delta\nu_{\mathrm{nt}}=w-\nu$，或者等价的，在式 (10.21) 中给出的 $\Delta\nu_{\mathrm{nt}}=\Delta\nu_{\mathrm{pd}}$。

\noindent\textbf{算法 10.2}\quad
不可行初始点 Newton 方法。

给定初始点 $x\in\operatorname{dom} f$，$\nu$，误差阈值 $\epsilon>0$，$\alpha\in(0,1/2)$，$\beta\in(0,1)$。

\noindent\textbf{重复进行}

\noindent 1. 计算原对偶 Newton 方向 $\Delta x_{\mathrm{nt}},\Delta\nu_{\mathrm{nt}}$。

\noindent 2. 对 $\|r\|_2$ 进行回溯直线搜索。

\[
t:=1.
\]
只要
\[
\|r(x+t\Delta x_{\mathrm{nt}},\nu+t\Delta\nu_{\mathrm{nt}})\|_2>(1-\alpha t)\|r(x,\nu)\|_2,
\]
令 $t:=\beta t$。

\noindent 3. 改进。$x:=x+t\Delta x_{\mathrm{nt}}$，$\nu:=\nu+t\Delta\nu_{\mathrm{nt}}$。

直到 $Ax=b$ 并且 $\|r(x,\nu)\|_2\leq \epsilon$。

该算法和处理可行初始点的标准 Newton 方法非常相似，仅有少数不同。首先，搜索方向包括额外的依赖原问题残差的修正项。其次，直线搜索目标是残差的范数，而不是函数 $f$。最后，算法停止时不仅要求原问题可行性被满足，还要求（对偶）残差的范数也很小。

对步骤 2 中的直线搜索有必要进行一些注释。同基于目标函数的直线搜索相比，对残差范数进行直线搜索可能增加成本，但其增量通常是微不足道的。同样，直线搜索一定在有限次迭代后结束，因为式 (10.23) 表明当 $t$ 很小时直线搜索的终止条件能够满足。

方程 (10.24) 表明，如果某次迭代的步长等于 $1$，下一个迭代点就是可行的，由此可知在那以后所有的迭代点都将是可行的。因此，一旦得到某个可行点，不可行初始点 Newton 方法给出的搜索方向，将和 \S10.2 描述的（可行）Newton 方法给出的搜索方向完全吻合。

关于不可行初始点 Newton 方法存在很多变形。例如，一旦满足了可行性，我们就可以切换到 \S10.2 描述的（可行）Newton 方法。（换言之，我们可以将直线搜索的目标函数换成 $f$，将终止条件换成 $\lambda(x)^2/2\leq \epsilon$。）当可行性满足时，不可行初始点 Newton 方法和标准（可行）Newton 方法只在回溯和终止条件方面有些差别，其算法性能非常相似。

\noindent\textbf{利用不可行初始点 Newton 方法简化初始化}\quad
不可行初始点 Newton 方法的主要优点在于对初始化的要求简单。如果 $\operatorname{dom} f=\mathbb{R}^n$，初始化（可行）Newton 方法就是计算 $Ax=b$ 的一个解，在这种情况下用不可行初始点 Newton 方法，除了方便一点，没有什么特别的优点。

如果 $\operatorname{dom} f$ 不等于 $\mathbb{R}^n$，在 $\operatorname{dom} f$ 找到满足 $Ax=b$ 的点本身就具有挑战性。一种一般性的处理方法，当 $\operatorname{dom} f$ 很复杂并且不知道是否和 $\{z\mid Az=b\}$ 相交时可能也是最好的方法，是采用阶段 I 方法（见 \S11.4）求出这样一个点（或者证实 $\operatorname{dom} f$ 和 $\{z\mid Az=b\}$ 不相交）。但如果 $\operatorname{dom} f$ 比较简单，并且已知含有满足 $Ax=b$ 的点，那么不可行初始点 Newton 方法就是一种简单的可选方案。

一种常见的例子发生于 $\operatorname{dom} f=\mathbb{R}_{++}^n$ 的情况，例如在例 10.2 中描述的等式约束解析中心点问题。为了对以下问题初始化 Newton 方法，
\begin{equation}
\begin{aligned}
\text{minimize}\quad & -\sum_{i=1}^n\log x_i\\
\text{subject to}\quad & Ax=b,
\end{aligned}
\tag{10.25}
\end{equation}
我们需要找到满足 $Ax=b$ 的 $x^{(0)}\succ 0$，这等价于求解一个标准形式的线性规划可行性问题。该问题可以采用阶段 I 方法处理，或者等价的，采用不可行初始点 Newton 方法，选取任何正分量的初始点，例如，$x^{(0)}=\mathbf{1}$。

同样的技巧可以用于求解无约束问题时未知 $\operatorname{dom} f$ 的初始点的情况。作为一个例子，我们考虑等式约束解析中心点问题 (10.25) 的对偶问题，
\[
\text{maximize}\quad g(\nu)=-b^T\nu+n+\sum_{i=1}^n\log(A^T\nu)_i.
\]
为了对该问题初始化（可行初始点）Newton 方法，我们必须找到满足 $A^T\nu^{(0)}\succ 0$ 的 $\nu^{(0)}$，即我们必须求解一组线性不等式。该问题可以用阶段 I 方法求解，或者在重新描述问题后用不可行初始点 Newton 方法求解。我们先采用新的变量 $y\in\mathbb{R}^n$ 将其表示为一个等式约束问题
\[
\begin{aligned}
\text{maximize}\quad & -b^T\nu+n+\sum_{i=1}^n\log y_i\\
\text{subject to}\quad & y=A^T\nu,
\end{aligned}
\]
然后就可以从任何正分量的 $y^{(0)}$（和任意的 $\nu^{(0)}$）开始应用不可行初始点 Newton 方法。

对于不知道是否存在严格可行的初始点的问题，采用不可行初始点 Newton 方法进行初始化存在一个缺点，这就是不能明确判断不存在严格可行点的情况；因为此时残差范数仍然会很慢地收敛于某个正数。（作为对比，阶段 I 方法能够清楚地判定这种情况。）此外，在可行性满足以前，不可行初始点 Newton 方法可能收敛地很慢；见 \S11.4.2。

\subsection*{10.3.3\quad 收敛性分析}

本节说明，如果一些假设成立，那么不可行初始点 Newton 方法将收敛到最优解。收敛性证明和标准 Newton 方法或者等式约束的标准 Newton 方法的收敛性证明非常相似。我们将证明，一旦残差范数足够小，算法将选取完全的步长（意味着可行性被满足），此后将保持二次收敛速度。我们也将证明，在二次收敛以前，每次迭代都能使残差范数至少减少一个固定的量。由于残差范数不能为负值，经过有限次迭代，残差必然减少到能够保证选取完全步长的程度，从而进入二次收敛。

\noindent\textbf{假设}\quad
我们提出以下假设。

\begin{itemize}
\item 下水平集
\begin{equation}
S=\{(x,\nu)\mid x\in\operatorname{dom} f,\ \|r(x,\nu)\|_2\leq \|r(x^{(0)},\nu^{(0)})\|_2\}
\tag{10.26}
\end{equation}
是闭的。如果 $f$ 是闭的，则 $\|r\|_2$ 是一个闭函数，此时该条件被任何 $x^{(0)}\in\operatorname{dom} f$ 和任何 $\nu^{(0)}\in\mathbb{R}^p$ 所满足（见习题 10.7）。

\item 在集合 $S$ 上对某些 $K$ 成立
\begin{equation}
\|Dr(x,\nu)^{-1}\|_2
=
\left\|
\begin{bmatrix}
\nabla^2 f(x) & A^T\\
A & 0
\end{bmatrix}^{-1}
\right\|_2
\leq K.
\tag{10.27}
\end{equation}

\item 对于 $(x,\nu),(\tilde{x},\tilde{\nu})\in S$，$Dr$ 满足 Lipschitz 条件
\[
\|Dr(x,\nu)-Dr(\tilde{x},\tilde{\nu})\|_2
\leq L\|(x,\nu)-(\tilde{x},\tilde{\nu})\|_2.
\]
（该式等价于 $\nabla^2 f(x)$ 满足 Lipschitz 条件；参见习题 10.7。）
\end{itemize}

下面将看到，这些假设意味着 $\operatorname{dom} f$ 和 $\{z\mid Az=b\}$ 相交，并存在一个最优解 $(x^\star,\nu^\star)$。

\noindent\textbf{和标准 Newton 方法比较}\quad
上述假设和 \S10.2.4（第 505 页）分析标准 Newton 方法时采用的假设非常相似。其中第二和第三个假设，即 KKT 的逆矩阵有界和 Lipschitz 条件，本质上相同。但是，对于不可行初始点 Newton 方法，下水平集条件 (10.26) 比 \S10.2.4 的下水平集条件更有一般性。

作为一个例子，考虑等式约束的最大熵问题
\[
\begin{aligned}
\text{minimize}\quad & f(x)=\sum_{i=1}^n x_i\log x_i\\
\text{subject to}\quad & Ax=b,
\end{aligned}
\]
其中 $\operatorname{dom} f=\mathbb{R}_{++}^n$。目标函数 $f$ 是非闭的；它有非闭的下水平集，因此，至少对于一些初始点，关于标准 Newton 方法的假定不成立。这里遇到的问题是，当 $x_i\to 0$ 时，负熵函数不趋于 $\infty$。但在另一方面，不可行初始点 Newton 方法的下水平集条件 (10.26) 对该问题能成立，这是因为当 $x_i\to 0$ 时负熵函数的梯度的范数趋于 $\infty$。因此，不可行初始点 Newton 方法可以保证解决等式约束的最大熵问题。（在这里我们并不知道标准 Newton 方法是否不能解决该问题；我们能够看到的仅是我们的收敛性分析不成立。）请注意，如果初始点满足等式约束，那么标准 Newton 方法和不可行初始点 Newton 方法的差别仅在于直线搜索，它只在阻尼阶段才会发生。

\noindent\textbf{一个基本不等式}\quad
我们从推导一个基本不等式开始。令 $y=(x,\nu)\in S$ 满足 $\|r(y)\|_2\neq 0$，令 $\Delta y_{\mathrm{nt}}=(\Delta x_{\mathrm{nt}},\Delta\nu_{\mathrm{nt}})$ 为 $y$ 处的 Newton 方向。定义
\[
t_{\max}=\inf\{t>0\mid y+t\Delta y_{\mathrm{nt}}\notin S\}.
\]
如果对所有的 $t\geq 0$ 成立 $y+t\Delta y_{\mathrm{nt}}\in S$，我们按照惯例定义 $t_{\max}=\infty$。否则，$t_{\max}$ 表示满足 $\|r(y+t\Delta y_{\mathrm{nt}})\|_2=\|r(y^{(0)})\|_2$ 的最小的正的 $t$。特别是，由此可知对任意的 $0\leq t\leq t_{\max}$ 都成立 $y+t\Delta y_{\mathrm{nt}}\in S$。

我们将证明对于任意的 $0\leq t\leq \min\{1,t_{\max}\}$ 都成立
\begin{equation}
\|r(y+t\Delta y_{\mathrm{nt}})\|_2
\leq (1-t)\|r(y)\|_2+(K^2L/2)t^2\|r(y)\|_2^2.
\tag{10.28}
\end{equation}
我们有
\[
\begin{aligned}
r(y+t\Delta y_{\mathrm{nt}})
&=r(y)+\int_0^1 Dr(y+\tau t\Delta y_{\mathrm{nt}})t\Delta y_{\mathrm{nt}}\,d\tau\\
&=r(y)+tDr(y)\Delta y_{\mathrm{nt}}
  +\int_0^1(Dr(y+\tau t\Delta y_{\mathrm{nt}})-Dr(y))t\Delta y_{\mathrm{nt}}\,d\tau\\
&=r(y)+tDr(y)\Delta y_{\mathrm{nt}}+e\\
&=(1-t)r(y)+e,
\end{aligned}
\]
其中利用了 $Dr(y)\Delta y_{\mathrm{nt}}=-r(y)$，并定义了
\[
e=\int_0^1(Dr(y+\tau t\Delta y_{\mathrm{nt}})-Dr(y))t\Delta y_{\mathrm{nt}}\,d\tau.
\]
现在假设 $0\leq t\leq t_{\max}$，因此对所有的 $0\leq\tau\leq 1$ 均成立 $y+\tau t\Delta y_{\mathrm{nt}}\in S$。我们可以对 $\|e\|_2$ 建立以下上界：
\[
\begin{aligned}
\|e\|_2
&\leq \|t\Delta y_{\mathrm{nt}}\|_2
\int_0^1\|Dr(y+\tau t\Delta y_{\mathrm{nt}})-Dr(y)\|_2\,d\tau\\
&\leq \|t\Delta y_{\mathrm{nt}}\|_2
\int_0^1 L\|\tau t\Delta y_{\mathrm{nt}}\|_2\,d\tau\\
&=(L/2)t^2\|\Delta y_{\mathrm{nt}}\|_2^2\\
&=(L/2)t^2\|Dr(y)^{-1}r(y)\|_2^2\\
&\leq (K^2L/2)t^2\|r(y)\|_2^2,
\end{aligned}
\]
其中第二行利用了 Lipschitz 条件，最后利用了上界 $\|Dr(y)^{-1}\|_2\leq K$。现在我们可以导出式 (10.28) 中的上界：对于任意的 $0\leq t\leq \min\{1,t_{\max}\}$，
\[
\begin{aligned}
\|r(y+t\Delta y_{\mathrm{nt}})\|_2
&=\|(1-t)r(y)+e\|_2\\
&\leq (1-t)\|r(y)\|_2+\|e\|_2\\
&\leq (1-t)\|r(y)\|_2+(K^2L/2)t^2\|r(y)\|_2^2.
\end{aligned}
\]

\noindent\textbf{阻尼 Newton 阶段}\quad
我们首先说明如果 $\|r(y)\|_2>1/(K^2L)$，不可行初始点 Newton 方法的一次迭代将把 $\|r\|_2$ 至少减少一个确定的量。

基本不等式 (10.28) 的右边是 $t$ 的二次函数，并在 $t=0$ 和它的最小解
\[
\bar{t}=\frac{1}{K^2L\|r(y)\|_2}<1
\]
之间单调减少。我们有 $t_{\max}>\bar{t}$，否则将可推出不正确的关系 $\|r(y+t_{\max}\Delta y_{\mathrm{nt}})\|_2<\|r(y)\|_2$。因此基本不等式在 $t=\bar{t}$ 时成立，于是
\[
\begin{aligned}
\|r(y+\bar{t}\Delta y_{\mathrm{nt}})\|_2
&\leq \|r(y)\|_2-1/(2K^2L)\\
&\leq \|r(y)\|_2-\alpha/(K^2L)\\
&=(1-\alpha\bar{t})\|r(y)\|_2,
\end{aligned}
\]
它表明步长 $\bar{t}$ 满足直线搜索的停止条件。因此我们有 $t\geq \beta\bar{t}$，其中 $t$ 是回溯算法确定的步长。从 $t\geq \beta\bar{t}$ 可得（根据回溯直线搜索算法的停止条件）
\[
\begin{aligned}
\|r(y+t\Delta y_{\mathrm{nt}})\|_2
&\leq (1-\alpha t)\|r(y)\|_2\\
&\leq (1-\alpha\beta\bar{t})\|r(y)\|_2\\
&=\left(1-\frac{\alpha\beta}{K^2L\|r(y)\|_2}\right)\|r(y)\|_2\\
&=\|r(y)\|_2-\frac{\alpha\beta}{K^2L}.
\end{aligned}
\]
于是，只要 $\|r(y)\|_2>1/(K^2L)$，每次迭代就能把 $\|r\|_2$ 至少减少 $\alpha\beta/(K^2L)$。由此可知，至多经过
\[
\frac{\|r(y^{(0)})\|_2K^2L}{\alpha\beta}
\]
次迭代就可得到 $\|r(y^{(k)})\|_2\leq 1/(K^2L)$。

\noindent\textbf{二次收敛阶段}\quad
现在假设 $\|r(y)\|_2\leq 1/(K^2L)$。基本不等式保证
\begin{equation}
\|r(y+t\Delta y_{\mathrm{nt}})\|_2
\leq (1-t+(1/2)t^2)\|r(y)\|_2
\tag{10.29}
\end{equation}
对任意的 $0\leq t\leq \min\{1,t_{\max}\}$ 均成立。此时我们有 $t_{\max}>1$，否则由式 (10.29) 可得 $\|r(y+t_{\max}\Delta y_{\mathrm{nt}})\|_2<\|r(y)\|_2$，它和 $t_{\max}$ 的定义相矛盾。因此不等式 (10.29) 对 $t=1$ 成立，即我们有
\[
\|r(y+\Delta y_{\mathrm{nt}})\|_2
\leq (1/2)\|r(y)\|_2
\leq (1-\alpha)\|r(y)\|_2.
\]
该式表明 $t=1$ 能够满足回溯直线搜索算法的停止条件，因此可以取完整步长。不仅如此，对于此后所有迭代我们都有 $\|r(y)\|_2\leq 1/(K^2L)$，于是在此后所有迭代中均可取完整步长。

我们可以把不等式 (10.28)（对于 $t=1$）写成
\[
\frac{K^2L\|r(y^+)\|_2}{2}
\leq
\left(\frac{K^2L\|r(y)\|_2}{2}\right)^2,
\]
其中 $y^+=y+\Delta y_{\mathrm{nt}}$。于是，如果用 $r(y^{+k})$ 代表 $\|r(y)\|_2\leq 1/K^2L$ 出现后再进行 $k$ 步迭代后的残差，我们有
\[
\frac{K^2L\|r(y^{+k})\|_2}{2}
\leq
\left(\frac{K^2L\|r(y)\|_2}{2}\right)^{2^k}
\leq
\left(\frac{1}{2}\right)^{2^k},
\]
即我们得到 $\|r(y)\|_2$ 相对零的二次收敛。

为了证明迭代序列收敛，我们将说明它是一个 Cauchy 序列。假设 $y$ 满足 $\|r(y)\|_2\leq 1/(K^2L)$，用 $y^{+k}$ 表示 $y$ 以后再进行 $k$ 次迭代得到的相应值。既然这些迭代均在二次收敛的区域内，步长都等于 $1$，因此我们有
\[
\begin{aligned}
\|y^{+k}-y\|_2
&\leq \|y^{+k}-y^{+(k-1)}\|_2+\cdots+\|y^+-y\|_2\\
&=\|Dr(y^{+(k-1)})^{-1}r(y^{+(k-1)})\|_2+\cdots+\|Dr(y)^{-1}r(y)\|_2\\
&\leq K\left(\|r(y^{+(k-1)})\|_2+\cdots+\|r(y)\|_2\right)\\
&\leq K\|r(y)\|_2\sum_{i=0}^{k-1}\left(\frac{K^2L\|r(y)\|_2}{2}\right)^{2^i-1}\\
&\leq K\|r(y)\|_2\sum_{i=0}^{k-1}\left(\frac{1}{2}\right)^{2^i-1}\\
&\leq 2K\|r(y)\|_2,
\end{aligned}
\]
其中第三行对所有迭代利用了假设 $\|Dr^{-1}\|_2\leq K$。因为 $\|r(y^{(k)})\|_2$ 收敛于零，我们可以断定 $y^{(k)}$ 是一个 Cauchy 序列，因此收敛。由 $r$ 的连续性可得，$y^\star$ 的极限满足 $r(y^\star)=0$。由此建立了我们较早的论断，即本节开始时的假设意味着存在最优解 $(x^\star,\nu^\star)$。

\subsection*{10.3.4\quad 凸-凹对策}

从不可行初始点 Newton 方法的收敛性证明可以发现，除了等式约束的凸优化问题，该方法还可用于更加广泛的一类问题。假设 $r:\mathbb{R}^n\to\mathbb{R}^n$ 是可微的，它的梯度在 $S$ 上满足 Lipschitz 条件，并且 $\|Dr(x)^{-1}\|_2$ 在 $S$ 上有界，其中
\[
S=\{x\in\operatorname{dom} r\mid \|r(x)\|_2\leq \|r(x^{(0)})\|_2\}
\]
是一个闭集。此时从任何 $x^{(0)}$ 开始不可行初始点 Newton 方法，都将收敛于 $r(x)=0$ 在 $S$ 中的一个解。在不可行初始点 Newton 方法中，我们将该性质应用于 $r$ 是等式约束优化问题残差的情况。然而它也可以应用于若干其他有意义的情况。一个有意义的例子就是求解凸-凹对策。（关于其他有关对策问题的讨论可参见 \S5.4.3 和习题 5.25）。

在 $\mathbb{R}^p\times\mathbb{R}^q$ 上的一个无约束（二人零和）对策由其支付函数 $f:\mathbb{R}^{p+q}\to\mathbb{R}$ 所定义。其含义是，参与对策的甲方选择一个值（或称为动作）$u\in\mathbb{R}^p$，参与对策的乙方选择一个值（或动作）$v\in\mathbb{R}^q$；基于这些选择，甲方支付乙方 $f(u,v)$。甲方的目的是极小化其支付量，而乙方的目的则是极大化甲方的支付量。

如果甲方首先选定 $u$，而乙方知道其选择值，那么乙方将选择使 $f(u,v)$ 达到最大的 $v$，由此产生支付量 $\sup_v f(u,v)$（假设上确界能够达到）。如果甲方预见到乙方将做出上述选择，他就应该选择使 $\sup_v f(u,v)$ 达到最小的 $u$。由此产生的甲方对乙方的支付量为
\begin{equation}
\inf_u\sup_v f(u,v),
\tag{10.30}
\end{equation}
（假设上确界能够达到。）另一方面，如果乙方首先做出选择，相反的推理成立，由此产生的甲方对乙方的支付量等于
\begin{equation}
\sup_v\inf_u f(u,v).
\tag{10.31}
\end{equation}
支付量 (10.30) 总是大于或等于支付量 (10.31)；两者之间的差距可以解释为知道对方的选择后再进行选择的优势。如果对所有的 $u,v$，
\[
f(u^\star,v)\leq f(u^\star,v^\star)\leq f(u,v^\star),
\]
我们就称 $(u^\star,v^\star)$ 是对策的一个解，或者是对策的一个鞍点。当对策有鞍点时，最后选择的优势将不复存在；$f(u^\star,v^\star)$ 将同时等于支付量 (10.30) 和 (10.31)。（参见习题 3.14。）

如果对于每个 $v$，$f(u,v)$ 是 $u$ 的凸函数，而对于每个 $u$，$f(u,v)$ 是 $v$ 的凹函数，相应对策被称为凸-凹对策。当 $f$ 可微（且凸-凹）时，对策的鞍点将满足方程 $\nabla f(u^\star,v^\star)=0$。

\noindent\textbf{用不可行初始点 Newton 方法求解}\quad
我们可以应用不可行初始点 Newton 方法计算具有二次可微支付函数的凸-凹对策问题的解。我们定义残差为
\[
r(u,v)=\nabla f(u,v)=
\begin{bmatrix}
\nabla_u f(u,v)\\
\nabla_v f(u,v)
\end{bmatrix},
\]
并对其应用不可行初始点 Newton 方法。对于对策问题，不可行初始点 Newton 方法被简单地称为 Newton 方法（针对凸-凹对策而言）。

只要 $Dr=\nabla^2 f$ 的逆矩阵有界，并在以下下水平集上满足 Lipschitz 条件：
\[
S=\{(u,v)\in\operatorname{dom} f\mid \|r(u,v)\|_2\leq \|r(u^{(0)},v^{(0)})\|_2\},
\]
其中 $u^{(0)},v^{(0)}$ 是对策双方的初始选择，我们就可以保证（不可行初始点）Newton 方法的收敛性。

和无约束优化问题中的强凸性条件有一个简单的类似之处，如果存在 $m>0$，对所有的 $(u,v)\in S$ 成立 $\nabla^2_{uu}f(u,v)\succeq mI$ 和 $\nabla^2_{vv}f(u,v)\preceq -mI$，我们就称支付函数 $f$ 是强凸-凹的。很自然地，这个强凸-凹性假设包含着逆矩阵有界假定（习题 10.10）。

\subsection*{10.3.5\quad 数例}

\noindent\textbf{一个简单的例子}\quad
我们用不可行初始点 Newton 方法求解等式约束的解析中心点问题 (10.25)。我们的第一个例子是随机产生的 $n=100$，$m=50$ 的实例，相应问题可行且存在下界。采用不可行初始点 Newton 方法，初始的原对偶变量取值为 $x^{(0)}=\mathbf{1}$，$\nu^{(0)}=0$，回溯直线搜索参数 $\alpha=0.01$，$\beta=0.5$。图 10.1 中的图像分别表示原对偶残差相对迭代次数的变化情况，而图 10.2 中的图像则显示对应的步长。第 8 次迭代时选取了完整的 Newton 步长，因此原残差变成（几乎等于）零，并在此后始终保持（几乎等于）零。在大约 9 次迭代后，（对偶）残差二次收敛于零。

\noindent\textbf{一个不可行的例子}\quad
我们再考虑一个和上面的例子维数相同的实例，不同之处在于 $\operatorname{dom} f$ 和 $\{z\mid Az=b\}$ 不相交，即这是一个不可行的问题。（该实例不满足问题 (10.1) 是可解的假设，也不满足 \S10.2.4 给出的假设；考虑该例仅仅为了说明当 $\operatorname{dom} f$ 和 $\{z\mid Az=b\}$ 不相交时应用不可行初始点 Newton 方法会发生什么情况。）图 10.3 给出了该实例的残差范数，图 10.4 给出了对应的步长。当然，在这种情况下步长永远不会等于 1，而残差也不会等于 0。

\noindent\textbf{一个凸-凹对策}\quad
我们最后考虑一个 $\mathbb{R}^{100}\times\mathbb{R}^{100}$ 上的凸-凹对策问题，支付函数
\begin{equation}
f(u,v)=u^TAv+b^Tu+c^Tv-\log(1-u^Tu)+\log(1-v^Tv),
\tag{10.32}
\end{equation}
其定义域
\[
\operatorname{dom} f=\{(u,v)\mid u^Tu<1,\ v^Tv<1\}.
\]
随机产生实例数据 $A,b$ 和 $c$。从 $u^{(0)}=v^{(0)}=0$ 开始（不可行初始点）Newton 方法，回溯参数选取 $\alpha=0.01$ 和 $\beta=0.5$，图 10.5 显示了迭代过程。

\section*{10.4\quad 实现}

\subsection*{10.4.1\quad 消除法}

为了实现消除法，我们必须计算满足下式的满秩矩阵 $F$ 和 $\hat{x}$
\[
\{x\mid Ax=b\}=\{Fz+\hat{x}\mid z\in\mathbb{R}^{n-p}\}.
\]
在 \S C.5 介绍的几种方法可以解决这个问题。

\subsection*{10.4.2\quad 求解 KKT 系统}

本节我们介绍计算 Newton 步径或不可行 Newton 步径的方法，其中都需要求解 KKT 形式的线性方程组
\begin{equation}
\begin{bmatrix}
H&A^T\\
A&0
\end{bmatrix}
\begin{bmatrix}
v\\
w
\end{bmatrix}
=-
\begin{bmatrix}
g\\
h
\end{bmatrix}.
\tag{10.33}
\end{equation}
此处我们假定 $H\in\mathbf{S}_+^n$，$A\in\mathbb{R}^{p\times n}$，并且 $\operatorname{rank}A=p<n$。类似方法可以用于计算凸-凹对策的 Newton 步径，其中系数矩阵的右下角块是负半定的（参见习题 10.13）。

\noindent\textbf{整体求解 KKT 系统}\quad
一个直接的方法就是简单求解由 $n+p$ 个变量的 $n+p$ 个线性方程组描述的 KKT 系统 (10.33)。KKT 矩阵是对称的，但可能不正定，因此一个好的做法是采用 LDL$^T$ 因式分解（参见 \S C.3.3）。如果不利用矩阵的结构，计算量是 $(1/3)(n+p)^3$ 次浮点运算。当问题较小时（即 $n$ 和 $p$ 都不大），或者 $A$ 和 $H$ 是稀疏矩阵时，这是一个合理的处理方法。

\noindent\textbf{通过消元求解 KKT 系统}\quad
通常比整体求解 KKT 系统更好的方法是基于变量 $v$ 的消元法（参见 \S C.4）。我们从描述最简单的情况开始，此时 $H\succ 0$。利用 KKT 方程组第一个方程
\[
Hv+A^Tw=-g,\qquad Av=-h,
\]
求解 $v$ 可得
\[
v=-H^{-1}(g+A^Tw).
\]
将其代入 KKT 方程组的第二个方程又可得到 $AH^{-1}(g+A^Tw)=h$，因此我们有
\[
w=(AH^{-1}A^T)^{-1}(h-AH^{-1}g).
\]
这些公式给了我们计算 $v$ 和 $w$ 的方法。

在计算 $w$ 的公式中出现的矩阵是 KKT 矩阵中 $H$ 的 Schur 补 $S$：
\[
S=-AH^{-1}A^T.
\]
由于 KKT 矩阵的特殊结构，加上我们已经假定 $A$ 的秩为 $p$，可知 $S$ 是负定的。

\noindent\textbf{算法 10.3\quad 采用消元法求解 KKT 系统}

给定 KKT 系统 $H\succ0$。

1. 计算 $H^{-1}A^T$ 和 $H^{-1}g$。

2. 计算 Schur 补 $S=-AH^{-1}A^T$。

3. 求解 $Sw=AH^{-1}g-h$ 确定 $w$。

4. 求解 $Hv=-A^Tw-g$ 确定 $v$。

步骤 1 可以采用先进行 $H$ 的 Cholesky 因式分解再递推计算 $p+1$ 个变量值的方式完成，计算量为 $f+(p+1)s$ 次浮点运算，其中 $f$ 是对 $H$ 进行因式分解的计算量，$s$ 为计算相关变量值的计算量。步骤 2 需要进行 $p\times n$ 和 $n\times p$ 的矩阵乘。如果不利用矩阵的结构，计算量为 $p^2n$ 次浮点运算。（因为结果是对称的，我们只需要计算 $S$ 的上三角部分。）在某些情况下可以利用 $A$ 和 $H$ 的结构更加有效地完成步骤 2。步骤 3 可以利用 $-S$ 的 Cholesky 因式分解完成，如果不进一步利用 $S$ 的结果，计算量为 $(1/3)p^3$ 次浮点运算。步骤 4 可以利用步骤 1 已经求出的 $H$ 的因式分解，因此计算量为 $2np+s$ 次浮点运算。如果在计算 Schur 补的过程中不利用问题的结构，那么总的浮点运算量为
\[
f+ps+p^2n+(1/3)p^3,
\]
（只保留主导项）。如果在计算 $S$ 时利用问题的结构，后两项可以更小。

如果能够对 $H$ 进行有效的因式化，那么采用分块消元方法比直接利用 LDL$^T$ 因式分解求解 KKT 系统在计算量上更有优势。例如，如果 $H$ 是对角阵（对应于可分目标函数的情况），我们有 $f=0$ 和 $s=n$，因此总的计算量是 $p^2n+(1/3)p^3$ 次浮点运算，它仅随 $n$ 线性增长。如果 $H$ 是带宽为 $k\ll n$ 的带状矩阵，那么 $f=nk^2$，$s=4nk$，因此总成本约为 $nk^2+4nkp+p^2n+(1/3)p^3$，仍然仅随 $n$ 线性增长。其他可以利用的 $H$ 的结构分别是对角块矩阵（对应于可分块的目标函数），稀疏矩阵，或对角加低秩矩阵；可参见附录 C 和 \S9.7 获得更多细节与例子。

\noindent\textbf{例 10.3\quad 等式约束的解析中心}\quad
我们考虑问题
\[
\begin{array}{ll}
\mbox{minimize} & -\displaystyle\sum_{i=1}^n\log x_i\\
\mbox{subject to} & Ax=b,
\end{array}
\]
此处目标函数是可分的，因此任意 $x$ 处的 Hessian 矩阵是对角阵
\[
H=\operatorname{diag}(x_1^{-2},\ldots,x_n^{-2}).
\]
如果我们采用一般方法，如利用 KKT 矩阵的 LDL$^T$ 因式分解，计算 Newton 方向，浮点运算成本为 $(1/3)(n+p)^3$。

如果采用分块消元方法计算 Newton 步径，成本为 $np^2+(1/3)p^3$，远小于一般方法的成本。

实际上这个成本和第 501 页例 10.2 中描述的计算对偶问题 Newton 步径的成本一样。对于（无约束）对偶问题，Hessian 矩阵是
\[
H_{\mathrm{dual}}=-ADA^T,
\]
其中 $D$ 是对角阵，其对角元素为 $D_{ii}=(A^T\nu)_i^{-2}$。计算这个矩阵的浮点运算成本为 $np^2$，而利用 $-H_{\mathrm{dual}}$ 的 Cholesky 因式分解确定 Newton 步径的成本为 $(1/3)p^3$。

\noindent\textbf{例 10.4\quad 等式约束的最小长度分片线性曲线}\quad
我们考虑 $\mathbb{R}^2$ 空间结点为 $(0,0),(1,x_1),\ldots,(n,x_n)$ 的分片线性曲线。为了确定满足等式约束 $Ax=b$ 的最小长度曲线，我们构造下述问题
\[
\begin{array}{ll}
\mbox{minimize} & (1+x_1^2)^{1/2}+\displaystyle\sum_{i=1}^{n-1}(1+(x_{i+1}-x_i)^2)^{1/2}\\
\mbox{subject to} & Ax=b,
\end{array}
\]
其中变量 $x\in\mathbb{R}^n$，$A\in\mathbb{R}^{p\times n}$。在这个问题中，目标函数是若干对相邻变量的函数和，因此 Hessian 矩阵 $H$ 是三对角阵。采用分块消元法计算 Newton 步径的浮点运算成本约为 $p^2n+(1/3)p^3$。

\noindent\textbf{$H$ 为奇异矩阵时的消元法}\quad
上面描述的分块消元方法显然不能用于 $H$ 是奇异矩阵的情况，但对原方法进行一些简单的改变就可以处理这种更一般的情况。这种更一般的方法是基于以下结果：KKT 是非奇异矩阵，当且仅当存在 $Q\succeq0$ 满足 $H+A^TQA\succ0$，在这种情况下对于所有的 $Q\succ0$ 成立 $H+A^TQA\succ0$。（见习题 10.1。）例如，我们可以推断，如果 KKT 矩阵非奇异，则 $H+A^TA\succ0$。

令 $Q\succeq0$ 是满足 $H+A^TQA\succ0$ 的一个矩阵。那么 KKT 系统 (10.33) 等价于
\[
\begin{bmatrix}
H+A^TQA&A^T\\
A&0
\end{bmatrix}
\begin{bmatrix}
v\\
w
\end{bmatrix}
=-
\begin{bmatrix}
g+A^TQh\\
h
\end{bmatrix}.
\]
由于 $H+A^TQA\succ0$，可以利用消元法求解该线性方程。

\subsection*{10.4.3\quad 例}

本节我们描述一些较长的例子，说明如何利用结构有效计算 Newton 步径。我们也给出一些数值结果。

\noindent\textbf{等式约束的解析中心}\quad
我们考虑等式约束的解析中心问题
\[
\begin{array}{ll}
\mbox{minimize} & f(x)=-\displaystyle\sum_{i=1}^n\log x_i\\
\mbox{subject to} & Ax=b.
\end{array}
\]
（见例 10.2 和例 10.3。）我们利用 $p=100$，$n=500$ 的一个问题比较三种方法。

第一种方法是等式约束的 Newton 方法（\S10.2）。Newton 步径 $\Delta x_{\mathrm{nt}}$ 由以下 KKT 系统 (10.11) 所定义：
\[
\begin{bmatrix}
H&A^T\\
A&0
\end{bmatrix}
\begin{bmatrix}
\Delta x_{\mathrm{nt}}\\
w
\end{bmatrix}
=
\begin{bmatrix}
-g\\
0
\end{bmatrix},
\]
其中 $H=\operatorname{diag}(1/x_1^2,\ldots,1/x_n^2)$，$g=-(1/x_1,\ldots,1/x_n)$。我们在第 522 页的例 10.3 中解释过，用消元法可以有效地求解 KKT 系统，即求解
\[
AH^{-1}A^Tw=-AH^{-1}g,
\]
并令 $\Delta x_{\mathrm{nt}}=-H^{-1}(A^Tw+g)$。换言之，
\[
\Delta x_{\mathrm{nt}}=-\operatorname{diag}(x)^2A^Tw+x,
\]
其中 $w$ 是以下方程的解，
\begin{equation}
A\operatorname{diag}(x)^2A^Tw=b.
\tag{10.34}
\end{equation}
图 10.6 给出误差和迭代次数的关系。不同曲线对应于四个不同的初始点。我们采用 $\alpha=0.1$，$\beta=0.5$ 的回溯直线搜索方法。

第二种方法是应用 Newton 方法求解对偶问题
\[
\mbox{maximize}\quad g(\nu)=-b^T\nu+\sum_{i=1}^n\log(A^T\nu)_i+n.
\]
（参见第 501 页的例 10.2。）此处 Newton 步径通过求解以下方程得到，
\begin{equation}
A\operatorname{diag}(y)^2A^T\Delta\nu_{\mathrm{nt}}=-b+Ay,
\tag{10.35}
\end{equation}
其中 $y=(1/(A^T\nu)_1,\ldots,1/(A^T\nu)_n)$。比较式 (10.35) 和式 (10.34) 我们看到两种方法具有相同的复杂性。图 10.7 给出了不同初始点产生的误差。我们采用了 $\alpha=0.1$，$\beta=0.5$ 的回溯直线搜索方法。

第三种方法是 \S10.3 介绍的不可行初始点 Newton 方法，我们将其应用于优化问题
\[
\nabla f(x^\star)+A^T\nu^\star=0,\qquad Ax^\star=b.
\]
Newton 步径通过求解以下方程获得
\[
\begin{bmatrix}
H&A^T\\
A&0
\end{bmatrix}
\begin{bmatrix}
\Delta x_{\mathrm{nt}}\\
\Delta\nu_{\mathrm{nt}}
\end{bmatrix}
=-
\begin{bmatrix}
g+A^T\nu\\
Ax-b
\end{bmatrix},
\]
其中 $H=\operatorname{diag}(1/x_1^2,\ldots,1/x_n^2)$，$g=-(1/x_1,\ldots,1/x_n)$。采用消元法可以有效求解该 KKT 系统，其计算成本和式 (10.34) 或式 (10.35) 相同。例如，如果首先求解
\[
A\operatorname{diag}(x)^2A^Tw=2Ax-b,
\]
则 $\Delta\nu_{\mathrm{nt}}$ 和 $\Delta x_{\mathrm{nt}}$ 由下式确定
\[
\Delta\nu_{\mathrm{nt}}=w-\nu,\qquad
\Delta x_{\mathrm{nt}}=x-\operatorname{diag}(x)^2A^Tw.
\]
图 10.8 显示对应于四个不同的初始点所产生的残差
\[
r(x,\nu)=(\nabla f(x)+A^T\nu,Ax-b)
\]
的范数和迭代次数的关系。我们采用了 $\alpha=0.1$，$\beta=0.5$ 的回溯直线搜索方法。

图形显示对该问题对偶方法似乎收敛速度最快，但也仅有 2 或 3 倍比例因子的差别。对偶方法大约经过 6 次迭代进入二次收敛，而原方法和不可行初始点 Newton 方法则分别经过 12--15 和 10--20 次迭代进入二次收敛。

这些方法对初始化的要求也不一样。原方法需要知道原问题的一个可行点，即满足 $Ax^{(0)}=b$，$x^{(0)}\succ0$ 的点。对偶方法需要对偶问题的一个可行解，即 $A^T\nu^{(0)}\succ0$。哪种初始点更容易得到，取决于具体问题。不可行初始点 Newton 方法不需要初始化；唯一的要求是 $x^{(0)}\succ0$。

\noindent\textbf{最优网络流}\quad
我们考虑由 $n$ 条边和 $p+1$ 个结点组成的连通有向图或网络。用 $x_j$ 代表通过边 $j$ 的流量或交通，$x_j>0$ 表示流量和边同向，$x_j<0$ 表示流量和边反向。对结点 $i$ 另外给定一个外部源流（或汇流）$s_i$，$s_i>0$ 表示进入结点，$s_i<0$ 表示离开结点。进出每个结点的流量必须满足流量守恒方程，其含义是进入一个结点的所有流量之和，包括外部源（汇）流，等于零。该守恒方程可以表示成 $\tilde A x=s$，其中 $\tilde A\in\mathbb{R}^{(p+1)\times n}$ 是有向图的结点进入矩阵，
\[
\tilde A_{ij}=
\begin{cases}
1 & \mbox{边 }j\mbox{ 离开结点 }i\\
-1 & \mbox{边 }j\mbox{ 进入结点 }i\\
0 & \mbox{其他情况。}
\end{cases}
\]
我们假定 $\mathbf{1}^Ts=0$，否则流量守恒方程 $\tilde A x=s$ 不能保持一致。（换言之，所有源流之和必须等于所有汇流之和。）由于 $\mathbf{1}^T\tilde A=0$，流量守恒方程 $\tilde A x=s$ 中存在多余的等式方程。我们可以消除任何一个方程得到独立的方程组，用 $Ax=b$ 表示，其中 $A\in\mathbb{R}^{p\times n}$ 是有向图的简化的结点进入矩阵（即删除一行的结点进入矩阵），而 $b\in\mathbb{R}^p$ 是简化的源向量（即删除相应分量的 $s$）。

总之，流量守恒由 $Ax=b$ 表示，其中 $A$ 是有向图的简化的结点进入矩阵，$b$ 是简化的源向量。矩阵 $A$ 是非常稀疏的矩阵，因为其每列至多只有两个非零元素（只能是 $+1$ 或 $-1$）。

我们给定所有源（汇）流，取交通流 $x$ 为变量，引入目标函数
\[
f(x)=\sum_{i=1}^n\phi_i(x_i),
\]
其中 $\phi_i:\mathbb{R}\to\mathbb{R}$ 表示边 $i$ 上的流量成本函数。我们假定流量成本函数是严格凸的二次可微函数。

在流量守恒约束下选择最佳流量的问题是
\begin{equation}
\begin{array}{ll}
\mbox{minimize} & \displaystyle\sum_{i=1}^n\phi_i(x_i)\\
\mbox{subject to} & Ax=b.
\end{array}
\tag{10.36}
\end{equation}
这里 Hessian 矩阵 $H$ 为对角矩阵，因为目标函数可分。

对于最优网络流问题 (10.36) 我们有多种计算 Newton 步径的方法。最直接的方法是利用稀疏的 LDL$^T$ 因式分解求解完全的 KKT 系统。

利用分块消除方法计算 Newton 步径可能是更好的选择。我们可以根据有向图确定 Schur 补 $S=-AH^{-1}A^T$ 的稀疏模式：当且仅当结点 $i$ 和结点 $j$ 由一个边相连时我们才有 $S_{ij}\neq0$。由此可知，如果网络是稀疏的，即每个结点由边连接到很少几个其他结点，那么 Schur 补 $S$ 是稀疏的。在这种情况下，我们可以利用稀疏性计算 $S$，进行有关的因式分解以及求解步骤。可以期望，计算 Newton 步径的计算复杂性随边的数目（即变量的数目）的增加而近似线性的增长。

\noindent\textbf{最优控制}\quad
我们考虑问题
\[
\begin{array}{ll}
\mbox{minimize} & \displaystyle\sum_{t=1}^N\phi_t(z(t))+\sum_{t=0}^{N-1}\psi_t(u(t))\\
\mbox{subject to} & z(t+1)=A_tz(t)+B_tu(t),\quad t=0,\ldots,N-1.
\end{array}
\]
此处

\begin{itemize}
\item $z(t)\in\mathbb{R}^k$ 是系统在 $t$ 时刻的状态，
\item $u(t)\in\mathbb{R}^l$ 是 $t$ 时刻对系统的控制，
\item $\phi_t:\mathbb{R}^k\to\mathbb{R}$ 是状态成本函数，
\item $\psi_t:\mathbb{R}^l\to\mathbb{R}$ 是输入成本函数，
\item $N$ 称为问题的时间视野。
\end{itemize}

假定输入和状态成本函数是严格凸的二次可微函数。问题变量为 $u(0),\ldots,u(N-1)$ 和 $z(1),\ldots,z(N)$。初始状态 $z(0)$ 给定。线性等式约束称为状态方程或动态演化方程。我们将整体优化变量 $x$ 定义为
\[
x=(u(0),z(1),u(1),\ldots,u(N-1),z(N))\in\mathbb{R}^{N(k+l)}.
\]
因为目标函数可分块（即一些 $z(t)$ 和 $u(t)$ 的函数之和），Hessian 矩阵是分块对角阵：
\[
H=\operatorname{diag}(R_0,Q_1,\ldots,R_{N-1},Q_N),
\]
其中
\[
R_t=\nabla^2\psi_t(u(t)),\quad t=0,\ldots,N-1,\qquad
Q_t=\nabla^2\phi_t(z(t)),\quad t=1,\ldots,N.
\]
我们将所有等式约束（即状态方程）写成 $Ax=b$，其中
\[
A=
\begin{bmatrix}
-B_0&I&0&0&0&\cdots&0&0&0\\
0&-A_1&-B_1&I&0&\cdots&0&0&0\\
0&0&0&-A_2&-B_2&\cdots&0&0&0\\
\vdots&\vdots&\vdots&\vdots&\vdots&\ddots&\vdots&\vdots&\vdots\\
0&0&0&0&0&\cdots&I&0&0\\
0&0&0&0&0&\cdots&-A_{N-1}&-B_{N-1}&I
\end{bmatrix},
\qquad
b=
\begin{bmatrix}
A_0z(0)\\
0\\
0\\
\vdots\\
0\\
0
\end{bmatrix}.
\]
$A$ 的行数（即等式约束数）是 $Nk$。采用稠密的 LDL$^T$ 因式分解直接求解 KKT 系统确定 Newton 步径，浮点运算量为
\[
(1/3)(2Nk+Nl)^3=(1/3)N^3(2k+l)^3.
\]

采用稀疏的 LDL$^T$ 因式分解可获得很大改进，因为这种方法能利用 $A$ 和 $H$ 中的大量零元素。

事实上，利用 $H$ 和 $A$ 的特殊的块状结构，采用分块消元方法计算 Newton 步径，可以取得更好的结果。此时 Schur 补 $S=-AH^{-1}A^T$ 具有 $k\times k$ 的块状三对角形式：
\[
S=-AH^{-1}A^T
=
\begin{bmatrix}
S_{11}&Q_1^{-1}A_1^T&0&\cdots&0&0\\
A_1Q_1^{-1}&S_{22}&Q_2^{-1}A_2^T&\cdots&0&0\\
0&A_2Q_2^{-1}&S_{33}&\cdots&0&0\\
\vdots&\vdots&\vdots&\ddots&\vdots&\vdots\\
0&0&0&\cdots&S_{N-1,N-1}&Q_{N-1}^{-1}A_{N-1}^T\\
0&0&0&\cdots&A_{N-1}Q_{N-1}^{-1}&S_{NN}
\end{bmatrix},
\]
其中
\[
S_{11}=-B_0R_0^{-1}B_0^T-Q_1^{-1},
\]
\[
S_{ii}=-A_{i-1}Q_{i-1}^{-1}A_{i-1}^T-B_{i-1}R_{i-1}^{-1}B_{i-1}^T-Q_i^{-1},\quad i=2,\ldots,N.
\]
特别是，由于 $S$ 是带宽为 $2k-1$ 的带状矩阵，我们可以用 $k^3N$ 的计算量对其因式分解。因此，当 $k\ll N$ 时，计算 Newton 步径的成本为 $k^3N$。该计算量随着时间视野 $N$ 的增加线性增长，而一般方法计算量的增长与 $N^3$ 同阶。

对于这个问题我们可以利用 $S$ 的分块三对角结构再前进一步。应用针对分块三对角矩阵的标准因式分解方法将导出求解二次最优控制问题的经典的 Riccati 递归。不过，如果只利用 $S$ 的带状性质仍然只能得到具有相同计算复杂性的算法。

\noindent\textbf{线性矩阵不等式的解析中心}\quad
我们考虑问题
\begin{equation}
\begin{array}{ll}
\mbox{minimize} & f(X)=-\log\det X\\
\mbox{subject to} & \operatorname{tr}(A_iX)=b_i,\quad i=1,\ldots,p,
\end{array}
\tag{10.37}
\end{equation}
其中 $X\in\mathbf{S}^n$ 是变量，$A_i\in\mathbf{S}^n$，$b_i\in\mathbb{R}$，$\operatorname{dom} f=\mathbf{S}_{++}^n$。该问题的 KKT 条件是
\begin{equation}
-X^{\star -1}+\sum_{i=1}^m\nu_i^\star A_i=0,\qquad
\operatorname{tr}(A_iX^\star)=b_i,\quad i=1,\ldots,p.
\tag{10.38}
\end{equation}
变量 $X$ 的维数是 $n(n+1)/2$。我们可以简单地忽略 $X$ 的矩阵结构，将其视为（向量）变量 $x\in\mathbb{R}^{n(n+1)/2}$，然后采用一般方法求解 $n(n+1)/2$ 个变量 $p$ 个等式约束的问题得到问题 (10.37) 的解。这样计算 Newton 步径的成本至少为
\[
(1/3)(n(n+1)/2+p)^3;
\]
它和 $n$ 的 $n^6$ 同阶。我们将看到存在一些更有吸引力的替代方案。

第一个选择是求解对偶问题。可求得 $f$ 的共轭
\[
f^\star(Y)=\log\det(-Y)^{-1}-n,
\]
而 $\operatorname{dom} f^\star=-\mathbf{S}_{++}^n$（见第 86 页的例 3.23），因此对偶问题为
\begin{equation}
\mbox{maximize}\quad -b^T\nu+\log\det\left(\sum_{i=1}^p\nu_iA_i\right)+n,
\tag{10.39}
\end{equation}
其定义域为
\[
\left\{\nu\mid \sum_{i=1}^p\nu_iA_i\succ0\right\}.
\]
这是变量为 $\nu\in\mathbb{R}^p$ 的无约束问题。最优的 $X^\star$ 可以利用最优的 $\nu^\star$ 通过求解式 (10.38) 中的第一个（对偶可行性）方程确定，即
\[
X^\star=\left(\sum_{i=1}^p\nu_i^\star A_i\right)^{-1}.
\]
考虑对偶问题 (10.39) 的 Newton 步径的计算成本。我们需要先计算 $g$ 的梯度和 Hessian 矩阵，然后确定 Newton 步径。梯度和 Hessian 矩阵由下式给出
\[
\nabla^2g(\nu)_{ij}=-\operatorname{tr}(A^{-1}A_iA^{-1}A_j),\quad i,j=1,\ldots,p,
\]
\[
\nabla g(\nu)_i=\operatorname{tr}(A^{-1}A_i)-b_i,\quad i=1,\ldots,p,
\]
其中 $A=\sum_{i=1}^p\nu_iA_i$。为了计算 $\nabla^2g(\nu)$ 和 $\nabla g(\nu)$，我们进行如下操作。首先，我们对每个 $j$ 计算 $A$（$pn^2$ 次浮点运算）和 $A^{-1}A_j$（$2pn^3$ 次浮点运算）。然后我们计算矩阵 $\nabla^2g(\nu)$。矩阵 $\nabla^2g(\nu)$ 有 $p(p+1)/2$ 个不同元素，每个元素都是 $\mathbf{S}^n$ 空间两个矩阵的内积，计算每个元素需要 $n(n+1)$ 次浮点运算，因此总的计算量（略去非主导项）为 $(1/2)p^2n^2$。既然已经有了 $A^{-1}A_i$，再计算 $\nabla g(\nu)$ 就只需很少计算量。最后，我们确定 Newton 步径 $-\nabla^2g(\nu)^{-1}\nabla g(\nu)$，成本为 $(1/3)p^3$ 次浮点运算。加在一起，并忽略非主导项，计算 Newton 步径的总成本为 $2pn^3+(1/2)p^2n^2+(1/3)p^3$ 次浮点运算。这是 $n$ 的 $n^3$ 函数，比上面描述的处理原问题的简单方法好很多，后者是 $n^6$ 函数。

我们也可以利用原问题的特殊的矩阵结构获得更加有效的求解方法。为了导出确定可行点 $X$ 处的 Newton 步径 $\Delta X_{\mathrm{nt}}$ 所需要的 KKT 系统，我们在 KKT 条件中用 $X+\Delta X_{\mathrm{nt}}$ 替换 $X^\star$，用 $w$ 替换 $\nu^\star$，并用一阶近似对第一个方程进行线性化
\[
(X+\Delta X_{\mathrm{nt}})^{-1}\approx X^{-1}-X^{-1}\Delta X_{\mathrm{nt}}X^{-1}.
\]
由此产生 KKT 系统
\begin{equation}
-X^{-1}+X^{-1}\Delta X_{\mathrm{nt}}X^{-1}
+\sum_{i=1}^p w_iA_i=0,\qquad
\operatorname{tr}(A_i\Delta X_{\mathrm{nt}})=0,\quad i=1,\ldots,p.
\tag{10.40}
\end{equation}
这是 $\Delta X_{\mathrm{nt}}\in\mathbf{S}^n$ 和 $w\in\mathbb{R}^p$ 的 $n(n+1)/2+p$ 个线性方程。如果用一般方法求解该方程组，计算成本的阶次为 $n^6$。

我们可以用分块消元方法更加有效地求解 KKT 系统 (10.40)。求解第一个方程可以确定 $\Delta X_{\mathrm{nt}}$，
\begin{equation}
\Delta X_{\mathrm{nt}}
=X-X\left(\sum_{i=1}^p w_iA_i\right)X
=X-\sum_{i=1}^p w_iXA_iX.
\tag{10.41}
\end{equation}
将该表达式代入另一个方程的 $\Delta X_{\mathrm{nt}}$ 可得
\[
\operatorname{tr}(A_j\Delta X_{\mathrm{nt}})
=\operatorname{tr}(A_jX)-\sum_{i=1}^p w_i\operatorname{tr}(A_jXA_iX)=0,\quad j=1,\ldots,p.
\]
这是 $w$ 的 $p$ 个线性方程：
\[
Cw=d,
\]
其中 $C_{ij}=\operatorname{tr}(A_iXA_jX)$，$d_i=\operatorname{tr}(A_iX)$。系数矩阵 $C$ 是对称正定阵，因此可以利用 Cholesky 因式分解确定 $w$。一旦有了 $w$，就可以用式 (10.41) 计算 $\Delta X_{\mathrm{nt}}$。

该方法的计算成本如下。首先计算 $A_iX$（$2pn^3$ 次浮点运算），然后计算矩阵 $C$。$C$ 有 $p(p+1)/2$ 个不同元素，每个是 $\mathbb{R}^{n\times n}$ 空间两个矩阵的内积，因此计算 $C$ 的成本为 $p^2n^2$ 次浮点运算。然后我们计算 $w=C^{-1}d$，其成本为 $(1/3)p^3$。最后我们计算 $\Delta X_{\mathrm{nt}}$。如果利用式 (10.41) 的第一个表达式，即先求和然后分别左乘和右乘 $X$，成本约为 $pn^2+3n^3$。所有加在一起，利用分块消元方法计算原问题的 Newton 步径的总成本为 $2pn^3+p^2n^2+(1/3)p^3$ 次浮点运算。这比阶次为 $n^6$ 的简单方法好很多，和计算对偶问题的 Newton 步径的成本一样。

\section*{参考文献}

我们在分析不可行 Newton 方法中采用的两个关键假设（导数 $Dr$ 存在有界逆和满足 Lipschitz 条件），对 Newton 方法的大多数收敛性证明都是核心条件；见 Ortega 和 Rheinboldt [OR00]以及 Dennis 和 Schnabel [DS96]。

在求解 KKT 系统时，对完整的系统直接进行因式分解，或采用消元法的相对好处已经在线性和二次规划的内点法框架中进行了广泛的研究；例如，见 Wright [Wri97, 第 11 章] 以及 Nocedal 和 Wright [NW99, \S16.1-2]。最优控制中的 Riccati 递推可以解释为对第 527 页的例子利用 $S$ 的 Schur 补的分块三对角结构的方法；该发现来自 Rao, Wright 和 Rawlings [RWR98, \S3.3]。

\section*{习题}

\noindent\textbf{等式约束优化}

\noindent\textbf{10.1 KKT 矩阵的非奇异性。}\quad
考虑 KKT 矩阵
\[
\begin{bmatrix}
P&A^T\\
A&0
\end{bmatrix},
\]
其中 $P\in\mathbf{S}_+^n$，$A\in\mathbb{R}^{p\times n}$，$\operatorname{rank}A=p<n$。

(a) 证明以下每个论断等价于 KKT 矩阵非奇异。
\begin{itemize}
\item $\mathcal{N}(P)\cap\mathcal{N}(A)=\{0\}$。
\item $Ax=0$，$x\neq0\Longrightarrow x^TPx>0$。
\item $F^TPF\succ0$，其中 $F\in\mathbb{R}^{n\times(n-p)}$ 是满足 $\mathcal{R}(F)=\mathcal{N}(A)$ 的矩阵。
\item 对某个 $Q\succeq0$ 成立 $P+A^TQA\succ0$。
\end{itemize}

(b) 证明如果 KKT 矩阵非奇异，它就正好有 $n$ 个正特征根和 $p$ 个负特征根。

\noindent\textbf{10.2 投影梯度法。}\quad
本题研究梯度法对等式约束优化问题的一种扩展。假设 $f$ 是凸的可微函数，$x\in\operatorname{dom} f$ 满足 $Ax=b$，其中 $A\in\mathbb{R}^{p\times n}$，$\operatorname{rank}A=p<n$。负梯度 $-\nabla f(x)$ 在 $\mathcal{N}(A)$ 上的 Euclid 投影由下式给出
\[
\Delta x_{\mathrm{pg}}=\mathop{\mathrm{argmin}}_{Au=0}\|-\nabla f(x)-u\|_2.
\]

(a) 令 $(v,w)$ 为以下方程的唯一解，
\[
\begin{bmatrix}
I&A^T\\
A&0
\end{bmatrix}
\begin{bmatrix}
v\\
w
\end{bmatrix}
=
\begin{bmatrix}
-\nabla f(x)\\
0
\end{bmatrix}.
\]
证明 $v=\Delta x_{\mathrm{pg}}$，$w=\mathop{\mathrm{argmin}}_y\|\nabla f(x)+A^Ty\|_2$。

(b) 假定 $F^TF=I$，那么投影负梯度 $\Delta x_{\mathrm{pg}}$ 和问题 (10.5) 的简化问题的负梯度之间有什么关系？

(c) 等式约束优化问题的投影梯度法采用步径 $\Delta x_{\mathrm{pg}}$，对 $f$ 进行回溯直线搜索。利用上面 (b) 的结果给出投影梯度法收敛于最优解的条件，假定从满足 $Ax^{(0)}=b$ 的 $x^{(0)}\in\operatorname{dom} f$ 开始迭代。

\noindent\textbf{10.3 对偶 Newton 方法。}\quad
本题研究用 Newton 方法求解等式约束优化问题 (10.1) 的对偶问题。假设 $f$ 二次可微，对所有 $x\in\operatorname{dom} f$ 成立 $\nabla^2 f(x)\succ0$，并且对每个 $\nu\in\mathbb{R}^p$，Lagrange 函数 $L(x,\nu)=f(x)+\nu^T(Ax-b)$ 有唯一最小解，用 $x(\nu)$ 表示。

(a) 证明对偶函数 $g$ 二次可微。对每个 $\nu$，用 $f$，$\nabla f$，以及在 $x=x(\nu)$ 处的 $\nabla^2 f$ 表示对偶函数 $g$ 的 Newton 步径。可以利用习题 3.40 的结果。

(b) 假设对所有的 $x\in\operatorname{dom} f$ 存在 $K$ 满足
\[
\left\|
\begin{bmatrix}
\nabla^2 f(x)&A^T\\
A&0
\end{bmatrix}^{-1}
\right\|_2\leq K.
\]
证明 $g$ 是强凹函数，满足 $\nabla^2 g(\nu)\preceq -(1/K)I$。

\noindent\textbf{10.4 简化问题的强凸性和 Lipschitz 常数。}\quad
假定 $f$ 满足 505 页给出的假设。证明简化的目标函数 $\tilde f(z)=f(Fz+\hat{x})$ 是强凸的，并且其 Hessian 矩阵是 Lipschitz 连续的（在相应的下水平集 $\tilde S$ 上）。用 $K$，$M$，$L$，以及 $F$ 的最大和最小奇异值表示强凸性和 $\tilde f$ 的 Lipschitz 常数。

\noindent\textbf{10.5 在目标中增加二次项。}\quad
假设 $Q\succeq0$。问题
\[
\begin{array}{ll}
\mbox{minimize} & f(x)+(Ax-b)^TQ(Ax-b)\\
\mbox{subject to} & Ax=b
\end{array}
\]
等价于原始的等式约束优化问题 (10.1)。该问题的 Newton 步径是否和原始问题的 Newton 步径相同？

\noindent\textbf{10.6 Newton 减量。}\quad
证明式 (10.13) 成立，即
\[
f(x)-\inf\{\hat f(x+v)\mid A(x+v)=b\}=\lambda(x)^2/2.
\]

\noindent\textbf{不可行初始点 Newton 方法}

\noindent\textbf{10.7 关于不可行初始点 Newton 方法的假设。}\quad
考虑 512 页给出的假设条件。

(a) 假设 $f$ 是闭函数。证明该假设意味着残差的范数 $\|r(x,\nu)\|_2$ 是闭函数。

(b) 证明 $Dr$ 满足 Lipschitz 条件的充要条件是 $\nabla^2 f$ 满足该条件。

\noindent\textbf{10.8 不可行初始点 Newton 方法和初始点已满足的等式约束。}\quad
假定在 $a_i^Tx=b_i$，$i=1,\ldots,p$ 的约束下用不可行初始点 Newton 方法极小化 $f(x)$。

(a) 假设初始点 $x^{(0)}$ 满足线性等式 $a_i^Tx=b_i$。证明以后的迭代点将始终满足该线性等式，即 $a_i^Tx^{(k)}=b_i$ 对所有 $k$ 均成立。

(b) 假定某个等式约束在第 $k$ 次迭代后被满足，即我们有 $a_i^Tx^{(k-1)}\neq b_i$，$a_i^Tx^{(k)}=b_i$。证明在第 $k$ 次迭代后，所有等式约束被满足。

\noindent\textbf{10.9 等式约束下的熵极大化。}\quad
考虑等式约束下的熵极大化问题
\begin{equation}
\begin{array}{ll}
\mbox{minimize} & f(x)=\displaystyle\sum_{i=1}^n x_i\log x_i\\
\mbox{subject to} & Ax=b
\end{array}
\tag{10.42}
\end{equation}
其中 $\operatorname{dom} f=\mathbb{R}_{++}^n$，$A\in\mathbb{R}^{p\times n}$。我们假设该问题是可行的，并且 $\operatorname{rank}A=p<n$。

(a) 证明该问题有唯一的最优解 $x^\star$。

(b) 确定 $A$，$b$ 和可行的 $x^{(0)}$ 使下水平集
\[
\{x\in\mathbb{R}_{++}^n\mid Ax=b,\ f(x)\leq f(x^{(0)})\}
\]
不是闭集。由此可知，某些可行初始点也不满足 505 页 \S10.2.4 列出的假设。

(c) 证明对于任何可行初始点，问题 (10.42) 满足 512 页 \S10.3.3 列出的关于不可行初始点 Newton 方法的假设。

(d) 导出 (10.42) 的 Lagrange 对偶，说明如何利用对偶问题最优解求得问题 (10.42) 的最优解。证明对于任意初始点，对偶问题满足 505 页 \S10.2.4 列出的假设。

以上 (b)，(c) 和 (d) 的结果并不意味着标准 Newton 方法失效，也不能说明不可行初始点 Newton 方法或者对偶方法在实践中工作得更好。它只表示当我们的收敛性分析可以应用于不可行初始点 Newton 方法和对偶方法时，却不能应用于标准 Newton 方法。（见习题 10.15。）

\noindent\textbf{10.10 强凸-凹对策的逆导数有界条件。}\quad
考虑支付函数为 $f$ 的凸-凹对策（见 517 页）。假设对任意的 $(u,v)\in\operatorname{dom} f$ 成立 $\nabla_{uu}^2 f(u,v)\succeq mI$ 和 $\nabla_{vv}^2 f(u,v)\preceq -mI$。证明
\[
\|Dr(u,v)^{-1}\|_2=\|\nabla^2 f(u,v)^{-1}\|_2\leq 1/m.
\]

\noindent\textbf{实现}

\noindent\textbf{10.11}\quad
考虑例 10.1 描述的资源配置问题。可以假设 $f_i$ 强凸，即对所有 $x$ 成立 $f_i''(x)\geq m>0$。

(a) 确定对简化问题计算 Newton 步径所需要的计算量。要求利用 Newton 方程的特殊结构。

(b) 说明如何通过对偶求解该问题。可以假设容易计算共轭函数 $f_i^\star$ 及其导数，并对给定的 $\nu$ 容易关于 $x$ 求解方程 $f_i'(x)=\nu$。给出计算对偶问题 Newton 步径的计算复杂性。

(c) 给出计算资源配置问题 Newton 步径的计算复杂性。要求利用 KKT 方程的特殊结构。

\noindent\textbf{10.12}\quad
给出计算下述问题 Newton 步径的一种有效方式，
\[
\begin{array}{ll}
\mbox{minimize} & \operatorname{tr}(X^{-1})\\
\mbox{subject to} & \operatorname{tr}(A_iX)=b_i,\quad i=1,\ldots,p,
\end{array}
\]
其定义域为 $\mathbf{S}_{++}^n$。假设 $p$ 和 $n$ 为阶次相同的数。另外导出其 Lagrange 对偶问题，并给出计算对偶问题 Newton 步径的计算复杂性。

\noindent\textbf{10.13 计算凸-凹对策 Newton 步径的消元法。}\quad
考虑支付函数为 $f:\mathbb{R}^p\times\mathbb{R}^q\to\mathbb{R}$ 的凸-凹对策（见 517 页）。假设 $f$ 是强凸-凹函数，即对所有 $(u,v)\in\operatorname{dom} f$ 和某个 $m>0$ 成立 $\nabla_{uu}^2 f(u,v)\succeq mI$ 和 $\nabla_{vv}^2 f(u,v)\preceq -mI$。

(a) 说明如何利用 $\nabla_{uu}^2f(u,v)$ 和 $-\nabla_{vv}^2f(u,v)$ 的 Cholesky 因式分解计算 Newton 步径。假设 $\nabla^2f(u,v)$ 稠密，比较以上方法和利用 $\nabla f(u,v)$ 的 LDL$^T$ 因式分解的方法的计算成本。

(b) 说明如何利用 $\nabla_{uu}^2f(u,v)$ 和/或 $\nabla_{vv}^2f(u,v)$ 的对角或分块对角结构。如果假设 $\nabla_{uv}^2f(u,v)$ 稠密，能够节省多少计算量？

\noindent\textbf{数值试验}

\noindent\textbf{10.14 对数-最优投资。}\quad
考虑习题 4.60 描述的对数-最优投资问题，没有 $x\succeq0$ 的约束。采用 Newton 方法计算以下问题数据对应的解：$n=3$ 种资产，$m=4$ 种情景，收益为
\[
p_1=
\begin{bmatrix}
2\\
1.3\\
1
\end{bmatrix},
\quad
p_2=
\begin{bmatrix}
2\\
0.5\\
1
\end{bmatrix},
\quad
p_3=
\begin{bmatrix}
0.5\\
1.3\\
1
\end{bmatrix},
\quad
p_4=
\begin{bmatrix}
0.5\\
0.5\\
1
\end{bmatrix}.
\]
四种情景的概率由 $\pi=(1/3,1/6,1/3,1/6)$ 给出。

\noindent\textbf{10.15 等式约束熵极大化。}\quad
考虑等式约束熵极大化问题
\[
\begin{array}{ll}
\mbox{minimize} & f(x)=\displaystyle\sum_{i=1}^n x_i\log x_i\\
\mbox{subject to} & Ax=b,
\end{array}
\]
其中 $\operatorname{dom} f=\mathbb{R}_{++}^n$，$A\in\mathbb{R}^{p\times n}$，$p<n$。（一些相关分析见习题 10.9。）生成一个 $n=100$，$p=30$ 的问题实例，随机选择 $A$（验证其为满秩阵），随机选择一个正向量作为 $\hat{x}$（例如，其分量在区间 $[0,1]$ 上均匀分布），然后令 $b=A\hat{x}$。（于是，$\hat{x}$ 可行。）

采用以下方法计算该问题的解。

(a) 标准 Newton 方法。可以选用初始点 $x^{(0)}=\hat{x}$。

(b) 不可行初始点 Newton 方法。可以选用初始点 $x^{(0)}=\hat{x}$（和标准 Newton 方法比较），也可以选用初始点 $x^{(0)}=\mathbf{1}$。

(c) 对偶 Newton 方法，即将标准 Newton 方法应用于对偶问题。

证实三种方法求得相同的最优点（和 Lagrange 乘子）。比较三种方法每步迭代的计算量，假设利用了相应的结构。（但在你的实现中不需要利用结构计算 Newton 步径。）

\noindent\textbf{10.16 凸-凹对策。}\quad
采用不可行初始点 Newton 方法求解形式如式 (10.32) 的凸-凹对策，随机产生数据。画出残差范数和步长关于迭代次数的图像。对直线搜索参数和初始点的影响进行实验（但要求它们必须满足 $\|u\|_2<1$，$\|v\|_2<1$）。

\end{document}
